% !TeX program = xelatex
% Build with XeLaTeX. Current implementation: MCM 1, RPX 11, RPA 12.
% RPX 11: exactly two proof blocks and fixed MI optimization equivalence.
\documentclass[UTF8,10pt,a4paper,oneside,fontset=fandol]{ctexbook}
\usepackage[top=25mm,bottom=25mm,left=26mm,right=24mm,headheight=15pt]{geometry}
\usepackage{amsmath,amssymb}
\usepackage{booktabs,longtable,tabularx,array}
\usepackage{xcolor,listings,fancyhdr,enumitem,microtype,etoolbox}
\usepackage{hyperref}
\usepackage{bookmark}
\definecolor{rpxblue}{HTML}{173A58}
\definecolor{rpxgray}{HTML}{F2F4F6}
\hypersetup{unicode=true,colorlinks=true,linkcolor=rpxblue,urlcolor=rpxblue,
  pdftitle={RPA/RPX Kernel 规范手册},pdfauthor={RPX kernel project},
  pdfsubject={机器模型、指令集、汇编、二进制格式与验证语义}}
\setlength{\parindent}{2em}
\setlength{\parskip}{3pt}
\setlength{\emergencystretch}{3em}
\setlist{nosep,leftmargin=2em,topsep=5pt}
\ctexset{chapter={format=\Large\bfseries\color{rpxblue},beforeskip=12pt,afterskip=20pt},
  section={format=\large\bfseries\color{rpxblue}},subsection={format=\normalsize\bfseries}}
\setcounter{tocdepth}{2}
\setcounter{secnumdepth}{2}
\pagestyle{fancy}
\fancyhf{}
\fancyhead[L]{\small RPA/RPX · Kernel 规范手册}
\fancyhead[R]{\small\nouppercase{\leftmark}}
\fancyfoot[L]{\small 2026-10-17 · MCM 1 · RPX 11 · RPA 12}
\fancyfoot[R]{\thepage}
\renewcommand{\headrulewidth}{0.4pt}
\fancypagestyle{plain}{\fancyhf{}\fancyfoot[L]{\small RPA/RPX · Kernel 规范手册}\fancyfoot[R]{\thepage}\renewcommand{\headrulewidth}{0pt}}
\lstset{basicstyle=\small\ttfamily,backgroundcolor=\color{rpxgray},
  frame=single,rulecolor=\color{black!15},framerule=0.3pt,framesep=5pt,
  breaklines=true,columns=fullflexible,keepspaces=true,showstringspaces=false,
  tabsize=4,aboveskip=10pt,belowskip=10pt}
% Keep each short listing and each complete example source on one page.
\BeforeBeginEnvironment{lstlisting}{\par\noindent\begin{minipage}{\linewidth}}
\AfterEndEnvironment{lstlisting}{\end{minipage}\par}
\newcommand{\rpxlisting}[2][]{\lstinputlisting[#1]{#2}}
\newcolumntype{Y}{>{\raggedright\arraybackslash}X}
\newcommand{\must}[1]{\textbf{必须：}#1}
% Inline code accepts raw underscores/dollars and escaped TeX punctuation.
\ExplSyntaxOn
\cs_generate_variant:Nn \tl_replace_all:Nnn {Nxx}
\cs_new_protected:Npn \rpx_code:n #1 {
  \group_begin:\ttfamily
  \tl_set:Nx \l_tmpa_tl {\tl_to_str:n {#1}}
  \tl_replace_all:Nxx \l_tmpa_tl {\tl_to_str:n {\%}} {\char_generate:nn {37}{12}}
  \tl_replace_all:Nxx \l_tmpa_tl {\tl_to_str:n {\_}} {\char_generate:nn {95}{12}}
  \tl_replace_all:Nxx \l_tmpa_tl {\tl_to_str:n {\$}} {\char_generate:nn {36}{12}}
  \tl_replace_all:Nxx \l_tmpa_tl {\tl_to_str:n {\#}} {\char_generate:nn {35}{12}}
  \tl_replace_all:Nxx \l_tmpa_tl {\tl_to_str:n {\&}} {\char_generate:nn {38}{12}}
  \tl_replace_all:Nxx \l_tmpa_tl {\tl_to_str:n {\{}} {\char_generate:nn {123}{12}}
  \tl_replace_all:Nxx \l_tmpa_tl {\tl_to_str:n {\}}} {\char_generate:nn {125}{12}}
  % Break long paths, signatures and camel-case names without adding glyphs.
  \regex_replace_all:nnN { ([a-z0-9])([A-Z]) } { \1\c{allowbreak}\2 } \l_tmpa_tl
  \regex_replace_all:nnN { ([/_.:,()<>]) } { \1\c{allowbreak} } \l_tmpa_tl
  \tl_use:N \l_tmpa_tl\group_end:}
\NewDocumentCommand{\code}{m}{
  \mode_if_math:TF {\text{\rpx_code:n {#1}}}{\rpx_code:n {#1}}}
\ExplSyntaxOff
\pdfstringdefDisableCommands{\def\code#1{#1}}
\begin{document}
\frontmatter
\hypersetup{pageanchor=false}
\pagenumbering{gobble}
\begin{titlepage}
\color{rpxblue}
\vspace*{12mm}
\noindent{\fontsize{54}{62}\selectfont\bfseries RPX}\par
\vspace{7mm}
\noindent{\Huge\bfseries RPA/RPX Kernel}\par
\vspace{4mm}
\noindent{\LARGE 精确重构证明规范手册}\par
\vspace{9mm}
\noindent\rule{\textwidth}{1pt}\par
\vspace{7mm}
\noindent{\large Reformulation Proof Assembly}\par
\noindent{\large Reformulation Proof eXecutable}\par
\vspace{10mm}
\noindent 完整模型 · 优化等价 · 推导与联合扩展\par
\noindent 文本汇编 · 二进制格式 · 完整证书\par
\noindent Loader · 执行语义 · 异常与符合性\par
\vfill
\noindent\textbf{当前实现参考 · MCM 1 · RPX 11 · RPA 12}\par
\noindent 2026 年 10 月 17 日\par
\vspace{5mm}
\noindent 依据当前安全 Rust kernel 的实现重写。\par
\noindent 证书仅含正向、反向两个完整证明块。\par
\end{titlepage}
\pagenumbering{roman}
\hypersetup{pageanchor=true}
\chapter*{阅读约定与文档状态}
\addcontentsline{toc}{chapter}{阅读约定与文档状态}

RPA 是 \emph{Reformulation Proof Assembly}；RPX 是 \emph{Reformulation Proof eXecutable}。
\code{.rpa} 是带版本头的严格 JSON 证书，\code{.rpx} 是同一证明对象的顺序二进制编码。
MCM 保存完整的精确有理系数模型。
调用者另行提供源模型、目标模型和所要求的 contract；证书不能自行降低这个要求。
本手册说明当前安全 Rust kernel 如何核验普通有理系数 MILP 的重构关系，
包括完整可行向量、目标关系、投影、双射、精确点见证及严格改进。

\section*{实现基线与协议状态}
\addcontentsline{toc}{section}{实现基线与协议状态}
本手册于 2026 年 10 月 13 日按实现重写，并于 10 月 14 日同步 V9/V10 冻结与统一 FORMAT 规范，取代 10 月 9 日的 D0 设计手册。
实现基线为 \code{crates/core/src/}：数学对象见 \code{model.rs} 与 \code{proof.rs}，
执行器见 \code{verify.rs}，文件编解码见 \code{codec.rs} 与 \code{rpa.rs}，
公开接口见 \code{api.rs}，命令行见 \code{main.rs}。
静态联合扩展分别由 flow、interval、menu、signed、subflow 模块实现；
point、injectivity、level 模块提供对应的精确 contract 检查。

\begin{longtable}{l p{104mm}}
\toprule 对象 & 当前状态 \\\midrule\endfirsthead
\toprule 对象 & 当前状态 \\\midrule\endhead
MCM & 只接受版本 1；与 RPX 的增量版本独立。\\
RPA / RPX 1--8 & 已实现、已冻结；后续版本沿用既有字段与数学意义。\\
RPA / RPX 9 & 已实现、已冻结的目标水平域；2026-10-14 冻结。\\
RPA / RPX 10 & 已实现、已冻结的严格改进、完整像排除、联合 subflow 与显式半空间域；2026-10-14 冻结。\\
D0 / v0.1 & 历史设计来源，既不是可接受的文本版本，也不是当前文件 ABI。\\
\bottomrule
\end{longtable}

实现完成与协议冻结分别登记。项目于 2026-10-14 明确冻结 V9/V10，
沿用既有字段、字节与数学接受条件；不兼容变更须使用新版本。
原 D0 的 section directory、表达式池、隐式 SSA 编号、逐项 mapping/check、
yield/merge 和可执行 theorem plugin 均不属于当前接口。
附录\ref{app:revisions}保留其设计决定与当前实现的对应关系。

\section*{规范性用语与证据}
\addcontentsline{toc}{section}{规范性用语与证据}
“必须”“不得”“拒绝”描述本版 kernel 的格式约束、数学接受条件和公开行为；
“建议”描述证书作者或集成方的工程选择；“可以”描述规则允许的情况。
数学语义、字段顺序、引用可见性和完整义务是规范性内容。
数据结构示意、性能解释、有限测试、历史测量和审阅记录是说明性证据。
一个文件通过汇编只说明能够编码，不等于证明通过验证。

\code{docs/FORMAT.md} 是 MCM 1 与 RPA/RPX 1--10 的唯一完整格式规范；
\code{audit/PROTOCOL_V9_V10_FREEZE.md} 登记冻结决定与验证身份。
各联合扩展的独立数学论证及审阅来源在相关章节和附录列出。
交付 manifest 固定的旧规范原字节保存在 \code{audit/specification-history/}；
历史验收记录保持交付时点的草案状态，不倒填冻结结论。
若发现手册与源代码不同，应将其作为需要修复的文档或实现问题，
不能让证书通过含糊解释、未定义行为或自动放宽检查取得接受结果。

\section*{可信边界}
\addcontentsline{toc}{section}{可信边界}
验证依赖实际完整模型、精确算术、静态推导规则和受审阅的联合存在引理。
可信实现包含 MPS/LP 导入器、MCM/RPX decoder、有理数库、表达式代数、
执行器及构造扩展模型的 builder。builder 的旧 scope 前提检查和模型合法性检查
并不能独立证明其生成的所有新坐标同时可行；该同时存在性质来自各自的可信引理。

RPA 作者、Python 生成器、求解器输出、审计 JSON、文件名和摘要都不能授予数学事实。
哈希用于记录输入来源；kernel 的接受条件仍然检查实际传入模型上的每一个前提和义务。
库的纯验证路径不读取文件系统、环境或时钟；文件操作与计时由 API 外层承担。
测试与独立有限 oracle 是实现证据，不能替代对 Rust 实现及依赖的形式化证明。

\section*{阅读顺序与符合性范围}
\addcontentsline{toc}{section}{阅读顺序与符合性范围}
数学审阅从第\ref{ch:model}章和第\ref{ch:execution}章开始：先确定所要求的命题，
再检查每个完整输出向量及分支、扩展、反证的语义。
第\ref{ch:instructions}章逐项列出推导与终结点的前置条件及拒绝条件；
第\ref{ch:format}章给出全部二进制记录，
第\ref{ch:loader}章说明真实输入方言、资源上限、API 和诊断，
第\ref{ch:assembly}章定义文本语法。
第\ref{ch:examples}章提供可重复运行的完整实例，
第\ref{ch:performance}章区分逻辑工作量、生成成本和进程测量。

\begin{description}
\item[Verifier] 接受所支持版本时，检查完整文件、独立 request、全部前提与实际模型义务；
不得把跳过的分支、失败的规则或资源耗尽报告为 \code{VERIFIED}。
\item[Assembler / disassembler] 负责严格文本与二进制的转换；
不承担独立数学认证，重新编码可以选择最低所需版本。
\item[证书作者] 安排已完成的推导与完整 mapping，保证每个可达叶使用同一个联合输出向量；
不能借用另一方向、兄弟分支或另一文件的局部事实。
\item[集成方] 明确传入完整模型、contract 和资源策略，依据 status 判断结果，
依据具体 contract 解释其数学结论。
\end{description}

本 kernel 不提供 proof search、一般 MILP 求解器、任意跳转、动态 theorem registry、
通用共享表达式 DAG、专用 batch 编码或数值解重建命令。
当前语料的完整验收不意味着任意重构都有紧凑证书；
\code{INVALID} 或 \code{RESOURCE_LIMIT} 也不意味着相应重构关系在数学上不成立。

\tableofcontents
\mainmatter
\chapter{数学模型与合约}
\label{ch:model}

本章定义验证器接受的数学对象，以及一次成功验证所证明的命题。
RPA 的含义为 \emph{Reformulation Proof Assembly}；RPX 的含义为
\emph{Reformulation Proof eXecutable}。
\code{.rpa} 是证明汇编源码，\code{.rpx} 是该源码的二进制证明证书。
RPX 的执行对象是数学事实、解映射和合约义务；它不执行求解器搜索，也不修改输入模型。

每个方向要证的是一组\textbf{验证义务}（obligation）：kernel 根据该方向的 output model
和所请求合约生成这些待证检查项（见第\ref{sec:directions}节与第\ref{sec:contracts}节），
证书必须逐条证明完毕；方向结束时仍为 \code{PENDING} 的义务以 \code{E_OBLIGATION} 拒绝。
每条义务由对象（输出变量的某个有限界或整数性、输出行的一个 side、目标 relation）
和所需事实类型（\code{LE}、\code{EQ} 或 \code{INT}）唯一确定，初始状态为 \code{PENDING}，
只有该方向 tape 中匹配的 \code{CHECK} 指令成功执行时才变为 \code{DONE}。
义务集合由 kernel 决定，证书无法选择其子集，也无法为义务自声明表达式；
待证表达式由 kernel 在检查时用当前 mapping 代换输出模型构造。

对证书而言，影响义务完成的途径只有三个：提供 mapping、提供从 input 静态事实和
分支假设出发的推导链、以及在 \code{CHECK} 中提交事实引用。
完成与否不取决于引用写法，而取决于提交事实的 $(\mathrm{kind},\mathrm{expression})$
是否与 kernel 构造的待证表达式精确相等；同一义务在不同分支可以对应不同的表达式。
分支内已完成的状态在汇合时按可达分支取交集。义务的生成规则与表达式构造见
第\ref{ch:execution}章第\ref{sec:obligation-generation}节，\code{CHECK} 指令族见
第\ref{ch:instructions}章，完整示例见第\ref{ch:examples}章。

\section{规范范围与模型边界}

v0.1 的数学域为具有有理数数据的普通混合整数线性规划（MILP）。
本规范不直接赋予浮点数、非线性表达式、半连续变量、半整数变量、SOS、
indicator constraint 或求解器专有对象任何语义。
这些对象若需使用本规范，必须先转换为本章定义的 MILP，且转换本身必须有明确的可信边界。

\code{MCM} 在本手册中表示 \emph{canonical mathematical model}，即规范化数学模型。
本章规定其数学 schema；它不规定 MCM 的文件字节布局。
\must{RPX 验证器必须针对调用者提供的、已经通过模型结构检查的规范化模型执行证明。}
证书中的模型摘要不得替代对模型结构的检查。

\paragraph{开放问题 O-01。}MCM 的文件格式、canonical serialization、SHA-256 输入字节以及
MPS 各种方言的导入规则尚未冻结。发布可互操作的 v0.1 工具链之前，必须另行规定它们。
在此之前，不得声称两个独立的 MCM 序列化器必然产生相同摘要。

\section{精确标量}

数学标量属于 $\mathbb{Q}$，并表示为
\[
  a=\frac{p}{q},\qquad p\in\mathbb{Z},\quad q\in\mathbb{Z}_{>0},
  \qquad \gcd(|p|,q)=1.
\]
零的唯一数值表示为 $0/1$。
负号属于分子；分母不得为零或负数。
表达式中的有限界、系数、目标常数和指令中的乘子均使用此数域。

\must{任何影响证明接受与否的运算和比较必须遵循精确有理数语义。}
实现可以采用机器整数快路径、大整数提升和经过证明正确的数值过滤，
但溢出、舍入或过滤不确定性不得改变数学结论。
IEEE 浮点数近似相等和容差可行性不属于本规范。

RPA 十进制和科学计数法字面量表示十进制有理数，例如
\[
  \text{\code{0.125}}=\frac18,\qquad
  \text{\code{1e-6}}=\frac{1}{1000000}.
\]
模型导入器必须说明原始文件数字的解释方式。
一个已经在外部系统中舍入的二进制浮点值，不能通过在证书中写出其十进制外观
而恢复舍入前的模型。

\section{规范化 MILP}
\label{sec:canonical-model}

一个模型写作
\[
  M=(V,R,f),\qquad V=(v_0,\ldots,v_{n-1}),\quad
  R=(r_0,\ldots,r_{m-1}).
\]
变量和行的编号是从零开始的稳定整数编号。
名称、注释、原始文件顺序以外的显示信息不参与任何数学推导。
模型中的每个引用必须在对应的编号范围内。

\subsection{变量}

变量 $v_i$ 定义一个实数坐标 $x_i$，并包含变量类型和有效上下界：
\[
  l_i\le x_i\le u_i,
  \qquad l_i\in\mathbb Q\cup\{-\infty\},\quad
  u_i\in\mathbb Q\cup\{+\infty\}.
\]
缺失的下界表示 $-\infty$，缺失的上界表示 $+\infty$。
无穷大是缺失界的标记，不是可参与有理数运算的标量。
每个变量的类型为以下一种：

\begin{center}
\begin{tabularx}{\textwidth}{lX}
\toprule
类型 & 附加数学条件 \\
\midrule
\code{CONTINUOUS} & 无附加整数条件，$x_i\in\mathbb R$。 \\
\code{INTEGER} & $x_i\in\mathbb Z$。 \\
\code{BINARY} & $x_i\in\mathbb Z$ 且 $0\le x_i\le1$。 \\
\bottomrule
\end{tabularx}
\end{center}

\must{BINARY 的隐含 $0$ 下界、$1$ 上界和整数性必须同时作为模型事实和
输出验证义务的一部分。}
若模型还提供额外二进制变量界，规范化阶段取交集，得到
$l_i=\max(0,l_i^{\rm explicit})$ 和
$u_i=\min(1,u_i^{\rm explicit})$。
模型 schema 要求有效界满足 $l_i\le u_i$。
本草案把违反此结构要求的模型视为无效模型；若需表示不可行模型，
可以通过合法的行与变量条件表示其矛盾。

\code{split.bin} 的前提是模型类型显式为 \code{BINARY}。
整数变量即使声明界为 $[0,1]$，也不能自动作为该指令的操作数；
它仍可使用整数分裂指令。

\subsection{约束行}

行 $r_j$ 的线性部分为
\[
  a_j(x)=\sum_{i\in I_j}a_{ji}x_i,
  \qquad a_{ji}\in\mathbb Q\setminus\{0\}.
\]
该行规定
\[
  L_j\le a_j(x)\le U_j,
  \qquad L_j\in\mathbb Q\cup\{-\infty\},\quad
  U_j\in\mathbb Q\cup\{+\infty\},\quad L_j\le U_j.
\]
没有有限界的行不产生可行性义务。
当 $L_j=U_j=b_j\in\mathbb Q$ 时，该行为等式行；否则为范围行。
单侧约束是范围行的特殊情形。
常数行合法，并可用于表达模型不可行性。

行的常数项统一移入界。
例如 $2x+3\le7$ 的模型行可表示为 $2x\le4$。
RPX fact 的表达式则保留常数项，不再单独存储右端项。
模型行和目标函数的非零系数列表同样必须按变量编号严格升序，
合并重复变量并删除零项；其有限标量遵循规范化有理数规则。

\subsection{目标函数}

规范化目标函数为
\[
  f_M(x)=c_M+\sum_i d_{Mi}x_i,
  \qquad c_M,d_{Mi}\in\mathbb Q,
\]
且目标方向统一为最小化。
原始最小化目标保持原值，原始最大化目标整体乘以 $-1$，包括目标常数。
原始方向和展示名称可以保留为外部元数据。

\proposal{P-01}{v0.1 统一使用最小化目标，并要求调用者明确要求的方向、
合约 profile 和 objective relation。证书声明必须与调用者请求一致；
不得由证书自行降低目标关系或省略一个请求的方向。}

这意味着目标偏移是证明的一部分。
删除或修改 $c_M$ 会改变目标关系，即使所有变量系数保持相同。
对于原始最大化问题，\code{LE} 关系在规范化目标上表示不增加，
在原始目标上对应不降低。

\subsection{可行域与扩展最优值}

模型的可行域定义为
\[
  \mathcal F(M)=
  \{x\in\mathbb R^n: \text{所有变量 domain 和所有约束行成立}\}.
\]
规范化最优值定义为扩展实数
\[
  \operatorname{val}(M)=\inf_{x\in\mathcal F(M)} f_M(x)
  \in\mathbb R\cup\{-\infty,+\infty\},
\]
其中 $\inf\varnothing=+\infty$。
非空可行域上的目标无下界时，最优值为 $-\infty$。
本定义不要求最优解存在。

\section{规范化仿射表达式}
\label{sec:canonical-affine}

仿射表达式写作
\[
  e(x)=c+\sum_{k=0}^{t-1}a_kx_{i_k}.
\]
其规范形式满足：
\begin{enumerate}
  \item $0\le i_0<i_1<\cdots<i_{t-1}<n$；
  \item 每个 $a_k$ 为非零的规范化有理数；
  \item 常数 $c$ 为规范化有理数；
  \item 同变量项已经合并，零系数项已经删除。
\end{enumerate}
空 term 列表表示常数表达式；常数为零时表示零表达式。
forward 和 reverse 各有独立的表达式池；池由有限个 entry 组成，每个 entry 承载一个上述规范形式的表达式，由 \code{ExprId} 标识，变量编号均指向该方向的 input model（编码见第\ref{ch:format}章）。

两个规范化表达式的相等是常数及各变量系数的精确数值相等。
不同 \code{ExprId} 可以指向数值相同的表达式。
不同 \code{ScalarRef} 也可能解码为同一个有理数。
\must{未经验证的编号相等、对象地址相等或散列值相等不得替代精确表达式比较。}
实现可以在经过验证的 interning 或精确比较结果上建立缓存。

规范化不包含按比例约简、改符号或寻找逻辑等价命题。
例如 $\mathrm{LE}(x-1)$ 与 $\mathrm{LE}(2x-2)$ 是等价不等式，
但在本 ISA 中是不同 proposition。
$\mathrm{EQ}(e)$ 与 $\mathrm{EQ}(-e)$ 同理。
若指令需要特定表达式，证明必须通过显式线性组合得到该表达式。

\section{事实类型}
\label{sec:fact-kinds}

事实（fact）由类型、表达式和有效作用域组成：
\[
  \mathrm{Fact}=(\mathrm{kind},\mathrm{expr},\mathrm{scope}).
\]
核心类型如下。

\begin{center}
\begin{tabularx}{\textwidth}{llX}
\toprule
类型 & 数学含义 & 示例 \\
\midrule
\code{LE} & $e(x)\le0$ & $\mathrm{LE}(3x_4-2x_9-7)$。 \\
\code{EQ} & $e(x)=0$ & $\mathrm{EQ}(x_3+x_7-1)$。 \\
\code{INT} & $e(x)\in\mathbb Z$ & $\mathrm{INT}(2x_3-x_8+5)$。 \\
\bottomrule
\end{tabularx}
\end{center}

fact 的 proposition 是二元组 $(\mathrm{kind},\mathrm{expr})$。
scope 限定该 proposition 在哪些分支条件下成立，但不参与 branch yield 的命题相等比较。
一个 fact 不携带独立 RHS、浮点容差、求解器 basis 或可执行代码。

\section{输入模型的静态事实}
\label{sec:static-model-facts}

以下引用直接从当前方向的 input model 获得事实，无需 \code{LOAD} 指令。
这里 $a_j$ 为行线性部分，$x_i$ 为输入变量。

\begin{center}
\begin{tabularx}{\textwidth}{l l X}
\toprule
RPA 引用 & FactKind & 表达式与前提 \\
\midrule
\code{@in.r[j].hi} & \code{LE} & $a_j(x)-U_j$；要求有限上界。 \\
\code{@in.r[j].lo} & \code{LE} & $-a_j(x)+L_j$；要求有限下界。 \\
\code{@in.r[j].eq} & \code{EQ} & $a_j(x)-b_j$；要求等式行 $L_j=U_j=b_j$。 \\
\code{@in.v[i].hi} & \code{LE} & $x_i-u_i$；要求有限有效上界。 \\
\code{@in.v[i].lo} & \code{LE} & $-x_i+l_i$；要求有限有效下界。 \\
\code{@in.v[i].int} & \code{INT} & $x_i$；要求 \code{INTEGER} 或 \code{BINARY}。 \\
\bottomrule
\end{tabularx}
\end{center}

等式行仍具有有限 hi 和 lo 事实；输出 obligation 则按等式行规则生成一个 EQ 义务。
连续变量不能提供 \code{.int} 事实，即使其有限上下界相同。
引用不存在的有限 side、错误类型的模型实体或越界编号必须拒绝。
所有静态模型事实在该方向的任意 scope 中有效。
分支假设 \code{@case} 属于另一类静态引用，见第\ref{ch:execution}章。

\section{两个执行方向}
\label{sec:directions}

一次 reformulation 验证绑定两个模型 $M_S$ 和 $M_T$。
每个启用的方向执行独立证明：

\begin{center}
\begin{tabularx}{\textwidth}{lllX}
\toprule
方向 & Input & Output & 要证明的解映射 \\
\midrule
\code{FORWARD} & $M_S$ & $M_T$ & $F:\mathcal F(M_S)\to\mathcal F(M_T)$。 \\
\code{REVERSE} & $M_T$ & $M_S$ & $G:\mathcal F(M_T)\to\mathcal F(M_S)$。 \\
\bottomrule
\end{tabularx}
\end{center}

fact register、表达式池、mapping、scope 栈和 obligation bitmap 均为 direction-local。
反向执行不得引用正向 fact 或依赖正向 mapping。
两方向共享证书的标量池和绑定的模型，但不共享证明状态。

一个方向的 mapping 必须在其每个可行输入上为每个输出变量提供唯一选定的值。
mapping 可以是单个仿射映射，也可以由有限的结构化分支表示。
验证器证明每个可达分支的输出都满足 output model；它不需要在验证时构造实际可行点。
重叠分支的结果可以不同，只要分别满足全部要求。
若需要输出可执行的点变换，必须同时规定确定的分支选择规则。

\section{目标关系与合约 profile}
\label{sec:contracts}

对一个方向，设 input model 为 $M_I$、output model 为 $M_O$、mapping 为 $E$。
目标差为
\[
  D_E(x)=f_{M_O}(E(x))-f_{M_I}(x).
\]
每个方向选择一种 objective relation：

\begin{center}
\begin{tabularx}{\textwidth}{l l X}
\toprule
关系 & 要求 & 验证义务 \\
\midrule
\code{NONE} & 无目标关系要求 & 不生成 objective obligation。 \\
\code{EQ} & $D_E(x)=0$ & 一个目标 EQ obligation。 \\
\code{LE} & $D_E(x)\le0$ & 一个目标 LE obligation。 \\
\bottomrule
\end{tabularx}
\end{center}

目标关系只在 input model 的可行域上要求成立。
不论 relation 为何，每个启用方向都必须完成全部 output mapping、变量 domain
和行可行性义务。
\code{NONE} 不允许略过模型可行性检查。

\subsection{方向性命题（资料性）}

数学上可以只要求一个独立方向。
例如一个 forward 证明意味着每个 source 可行点都有满足指定目标关系的 target 可行点。
它不能单独声称两个模型可行性或最优值等价。
本草案基础格式不编码这种单方向合约；它是未来扩展可采用的数学命题。
当前合法 RPX 必须同时启用两个方向，并选择下面两个等价 profile 中的恰好一个。

\subsection{可行性等价}

\code{feasibility-equivalent} profile 要求 forward 和 reverse 同时启用。
objective relation 可以分别为 \code{NONE}、\code{EQ} 或 \code{LE}。
成功验证证明
\[
  \mathcal F(M_S)=\varnothing
  \quad\Longleftrightarrow\quad
  \mathcal F(M_T)=\varnothing.
\]
它还证明所给的两个可行解转换存在；这不意味着两个不同维度可行域逐点相同。
若两个方向均选择 \code{NONE}，不声称目标值等价。

\subsection{最优值等价}

\code{opt-equivalent} profile 要求 forward 和 reverse 同时启用，
且每个方向的 objective relation 必须为 \code{EQ} 或 \code{LE}。
由两个方向分别得到
\[
  f_T(F(x))\le f_S(x),\qquad
  f_S(G(y))\le f_T(y),
\]
其中 EQ 是 LE 的特殊情形。
对两个不等式分别取下确界，得到
\[
  \boxed{\operatorname{val}(M_S)=\operatorname{val}(M_T)}.
\]
该结论适用于空可行域和目标无下界的情形。
映射同时保证两个可行域非空性一致，因而不会把不可行性与目标无下界混淆。
若两者具有有限最优值，任一侧的最优可行点通过相应映射成为另一侧的最优可行点。

\must{最优值等价不要求 $G\circ F$ 或 $F\circ G$ 为恒等映射。}
变量删除、冗余变量引入和支配关系变换通常没有双射解映射。
\code{opt-equivalent} 也不保证每个次优解的目标值完全保持不变；
该更强要求应明确选择两个方向的 \code{EQ}。

\subsection{合约标志的约束}

二进制 contract flags 的字段布局见二进制格式章。
本章规定它们的语义有效性：
\begin{enumerate}
  \item forward 和 reverse 都必须启用；
  \item 两个等价 profile 必须恰好声明一个，不得同时声明或同时省略；
  \item optimum profile 要求两个方向的 objective relation 均非 \code{NONE}；
  \item reserved relation 和 reserved flag 位必须为零，或按格式章规定拒绝。
\end{enumerate}
调用者请求与证书声明不一致属于 \code{E_CONTRACT}，
不允许通过读取证书后自动调整请求来消除该错误。

\section{摘要绑定与接受结果的含义}

RPX header 包含 source 和 target 的 SHA-256 摘要。
验证器必须比较它们与调用者提供的两个 canonical model 的摘要；
方向交换不能交换 header 的 source/target 字段含义。
模型摘要不匹配必须拒绝为 \code{E_HASH}。

摘要只有在明确的 canonical serialization 约定下才有可互操作意义。
调用者、MCM 导入器和验证器必须采用同一约定，并明确哪些外部导入步骤属于可信计算基。
RPX 成功证明的是绑定后两个数学模型之间的合约。
它不自动证明外部 MPS 文件与 MCM 完全一致，也不证明模型名、原始浮点缓冲区
或求解器中隐藏的数据与已绑定模型相同。

完整接受结果应至少报告：绑定的两个摘要、合约 profile、启用方向和各方向的目标关系。
只有所有启用方向成功，整个 RPX 才能报告 \code{VERIFIED}。

\chapter{执行语义与验收}
\label{ch:execution}

Certificate 恰含 forward 与 reverse。每个 Block 先执行有序 Steps，再执行一个 End。
两个根均从各自完整模型、空局部事实和空路径假设开始，共享同一个工作预算。
第一方向成功不会给第二方向增加事实，任一方向失败均不接受整个证书。
Block 证明完整输出构造；Derivation 只为 Cases 证明一条共同事实。

\section{引用解析与局部定义}
\label{sec:fact-numbering}

九种 Ref 的完整语义如下。编号均是相应逻辑空间的索引，而非文件偏移。
\begin{longtable}{p{0.22\textwidth}p{0.70\textwidth}}
\toprule
引用 & 解析结果与有效条件 \\
\midrule
\endhead
\code{VarLower}(i) & 当前 $I$ 的 $\mathrm{LE}(l_i-x_i)$；须存在有限下界。\\
\code{VarUpper}(i) & 当前 $I$ 的 $\mathrm{LE}(x_i-u_i)$；须存在有限上界。\\
\code{VarInt}(i) & 当前 $I$ 的 $\mathrm{INT}(x_i)$；须声明整数。\\
\code{RowLower}(j) & $\mathrm{LE}(L_j-a_j)$；须有有限下侧。\\
\code{RowUpper}(j) & $\mathrm{LE}(a_j-U_j)$；须有有限上侧。\\
\code{RowEq}(j) & $\mathrm{EQ}(a_j-L_j)$；须有相等的两有限侧。\\
\code{Local}(id) & $\Delta$ 中当前可见的已验定义；已 Drop、未定义或仅在兄弟可见时拒绝。\\
\code{Assumption}(k) & $\Gamma_k$；0 为最早加入的假设，越界拒绝。\\
\code{Exact} & 只对 checker 已确定的目标 proposition 做窄规则检查，不解析为普通事实。\\
\bottomrule
\end{longtable}

Lin、Round、EqPair、Cases、Plotkin 用显式 u32 \code{id} 发布一个局部事实。
ID 不须从零连续编号，也不须递增。
发布时该 ID 不得仍然活跃，且活跃数不得超过预算；所有前提、类型和表达式先通过检查。
Drop 按其数组顺序删除当前可见 ID，缺失 ID 或同一数组中重复删除已删 ID 均拒绝。
一个 ID 仅在删除后才能复用，新的定义有新的 proposition；引用只读当前定义。

\subsection{父子作用域}
\label{sec:branch-transactions}

Split 和 Cases 的两个子作用域继承父局部事实与既有假设。
子体可删除继承 ID，在该子体内复用 ID，并创建其他事实；这些变化不能改变父或兄弟。
两子体各附加本分支假设，退出时撤去该假设。
Cases 两子体都成功后，仅在父作用域发布其指定的一条共同结论；
子体的其他事实、删除或重定义均不向父发布。
Split 是端点，不存在需要合并的父续体。

lift 把当前活跃事实和路径假设带入其唯一子体。
旧坐标和旧行的完整前缀不变，因而旧事实仍有效；所有新行只在子体可见。
在分支内扩展不会给兄弟增加坐标或行。
实现可共享不可变父状态并保存子差量；这种存储选择不能改变上述可见性。

\section{目标命题检查与 Exact}
\label{sec:exact-check}

设 checker 要证明 $(K,e)$，通常引用 $r$ 解析为 $(K_r,e_r)$。
接受当且仅当 $e_r=e$ 精确成立，且 $K_r=K$ 或 $K_r=\code{Eq}$。
比例相等、仅在某个点取值相等、ID 相等及显示文本相等均不足。

Exact 不从证书读入命题；checker 先构造 $e$，再检查：
\begin{enumerate}
 \item EQ：$e$ 无变量项且常数为零。
 \item LE：$e$ 无变量项且常数不正。
 \item INT：常数为整数，每项系数为整数，且每个出现坐标声明整数。
\end{enumerate}
常数整数没有坐标前提。
非恒真的仿射不等式不能通过 Exact；无模型整数性支持的非恒定表达式也不能通过 Exact。
生成的菜单或 signed 二元坐标可满足 INT 的窄规则，
生成的 flow、interval 和 subflow 连续坐标不能。

Exact 可出现在 Return 检查、lift 前提和 DerivationEnd Witness。
它不能作为 Lin（即使乘子为零）、Round、EqPair、Plotkin、整数 Split/Cases
或任何 Absurd 的原始事实操作数。
一个已有 Eq 可做目标检查的种类弱化，但不赋予原始操作数自动种类转换。

\section{步骤序列与穷尽分类}

步骤顺序执行；后一步仅能引用当前有效的模型事实、假设与前面活跃定义。
各条推理规则见第\ref{ch:instructions}章。
六种 Step 的命名空间与八种 End、两种 DerivationEnd 的命名空间独立。
同一个数字 tag 在不同位置不表示同一操作。

Split 与 Cases 都先验证规范表达式 $e$。
未提交整数引用时，两个分支假设为
\[
 \mathrm{LE}(e),\qquad\mathrm{LE}(-e).
\]
它们覆盖全实数域，$e=0$ 的边界同时属于两支。
提交整数引用时，须解析出\emph{种类确为 Int}且表达式恰为 $e$ 的事实，
然后两假设为
\[
 \mathrm{LE}(e),\qquad\mathrm{LE}(1-e).
\]
整数域中分别表示 $e\le0$ 与 $e\ge1$，穷尽且互斥。
不能用 Exact 或 Eq 原始引用冒充该 Int 前提。

\subsection{Cases 的共同事实}

Cases 在父作用域指定 ID、Kind 和 expression；两 Derivation
分别在穷尽分支上证明这同一 proposition，或用普通 Absurd 证明该支为空。
Witness 使用目标命题检查规则，因此允许 Exact 以及同表达式的 Eq 种类弱化。
两支不能各自选择不同 expression，再要求 checker 猜测共同结论。
Derivation 没有 map、Return、lift。
内部 Steps 仍可递归 Cases；成功后父得到唯一指定事实。

\subsection{Split 的完整子证明}

每个 Split 子 Block 都必须以该上下文允许的端点成功完成。
两支 Return 可以给出不同完整输出向量；Absurd 可关闭已证明不可能的支。
不能把一个分支只验了变量界、另一分支只验了行的证明合成一个 Return。
判为不可能的支仍须先证明合法矛盾；kernel 不按猜测的可达性跳过证明。


\section{完整 Return 与义务生成}
\label{sec:mapping-state}
\label{sec:obligation-generation}
map 长度恰等于输出变量数，表达式使用当前输入坐标并须规范。
对同一个完整向量 $Y$，checks 依次证明：
\begin{enumerate}
\item 输出变量按坐标序：先有限下界 $l_j-Y_j\le0$，再有限上界 $Y_j-u_j\le0$，
然后对每个实际整数变量证明 $Y_j\in\mathbb Z$。
\item 输出行按行序：两有限侧相等时证明 $a_j[Y]-L_j=0$；否则先证明有限下侧，
再证明有限上侧。不存在的界不产生义务，矛盾界仍须检查。
\item 最后一项永远为 $f_O[Y]-f_I\le0$，包含两个完整目标的常数。
\end{enumerate}
每项引用必须证明准确表达式与种类，缺失、多余或错误顺序均拒绝。
等式事实可以证明同表达式的不增目标；没有 Equal/None 模式。
每个叶子的一套义务始终使用一个共同 $Y$。

\section{lift 的执行协议}
\label{sec:lift-execution}
五种 lift 均先检查 witness、尺寸与资源，再在\emph{扩展前}作用域核验全部前提，
最后进入生成模型上的唯一 body。旧模型的变量、行和目标保持完整前缀，新变量目标系数为零。
新事实只在 body 可用，不能循环证明本次 lift 的前提。嵌套 lift 可用上层已经建立的坐标。
存在引理保证一个共同扩展点满足全部新条件；内核构造公式，不运行数值选择或路径搜索。

\section{完成、资源和拒绝}
完整解码先于执行，所有分支字节均须合法。普通 Absurd 只接受正的恒定 Le 或非零恒定 Eq。
每个 Block 可以 Return、Split、Absurd 或五种 lift 完成；Cases 的 Derivation 只能 Witness 或 Absurd。
没有附加性质阶段，没有可跳过的方向，也不以猜测不可达为由略过失败分支。

默认工作预算为 $10^9$、活跃事实预算为 $10^6$、组合嵌套深度为 64。
两个方向共享工作预算，局部事实分别重建。Split、Cases 和 lift body 增加深度。
超出工作、深度、位数、项数或分配限制均不接受；资源不足不是数学反例。
诊断只报告 direction、step、obligation 等已知位置，不参与决定接受。

\section{正确性不变式}
\label{sec:soundness-invariants}
对完整输入域与路径假设归纳：可见事实有效；每个分类穷尽；矛盾叶没有可行点；
每个 lift 对每个父点给出共同可行扩展；每个 Return 证明完整输出可行且目标不增。
因此每个成功方向给出一个总构造，两个方向共同推出优化等价。
这依赖精确算术、模型导入、作用域、义务生成与可信构造规则的实现正确性，
不声称 Rust 程序已由证明助手形式化验证。

\chapter{核心指令集参考}
\label{ch:instructions}

本章规定 RPA/RPX 草案版本 0.1 的 28 条核心指令。每个条目同时给出文本汇编形式、二进制载荷、前置条件、数学推导规则、状态变化与拒绝条件。指令集作用于当前执行方向的输入模型与输出模型；正向和反向使用相同操作码，不共享寄存器、映射或义务状态。

\section{公共编码与执行约定}

\subsection{指令头与字节序}

所有指令采用以下 8 字节头。多字节字段均以 little-endian 编码；字段没有 C ABI 隐式填充。

\begin{lstlisting}
InstHeader {
    opcode:        u8;
    flags:         u8;
    aux:           u16;
    payload_bytes: u32;
}
\end{lstlisting}

指令地址必须 8 字节对齐。令当前地址为 \(p\)，则本指令占用的完整区间为
\[
 [p,\operatorname{align}_8(p+8+\mathrm{payload\_bytes})).
\]
载荷结束与下条指令之间的填充字节必须为零，并且不计入 \code{payload_bytes}。本章列出的载荷长度为精确长度；额外字段、截断字段或未解释的尾字节均以 \code{E_FORMAT} 拒绝。所有未使用的 \code{flags}、\code{aux} 与显式 \code{reserved} 字段必须为零。对长度、计数、偏移的检查必须使用无溢出的算术。

\subsection{操作数与寄存器定义}

\code{ExprId}、\code{FactRef}、\code{ScalarRef}、输入及输出对象索引都是 32 位字段。表达式编号仅在当前方向的表达式池中有效。\code{FactRef} 可以指向活跃派生事实、输入模型声明的事实，或当前可见分支假设；取值规则见机器模型章。

事实产生指令没有显式目标寄存器字段。其 \code{out_expr} 仅标识所验证的表达式。汇编中的目标寄存器由整个方向的隐式 SSA 编号规则确定，不能改变二进制执行顺序。\code{SPLIT} 的分支局部定义先占用编号，随后为汇合输出分配编号；编号细节与不可达代码的处理见机器模型章。所有生成的事实均具有当前 scope。

下文以 \(\operatorname{expr}(r)\) 和 \(\operatorname{kind}(r)\) 表示事实引用的表达式与类型。表达式相等是精确有理数意义下逐项相等：常数项、变量编号和对应系数必须相同。它不是 ExprId 相同，也不允许把正比例表达式视为同一命题。

\subsection{校验与提交}

每条指令必须先完成载荷、引用、类型与数学前置条件检查，再提交可观察的状态变化。失败不能提交部分事实、部分映射或部分义务。\code{E_REF} 表示引用不存在、已死亡、来自未来或不可见；\code{E_KIND} 表示事实类型不符合该指令；\code{E_PROOF} 表示所声明的推导不成立。错误符号的完整定义见异常章。

执行只计算精确有理数。任何实现都不得用浮点近似决定表达式相等、符号、整数性、取整或分支覆盖。宿主整数溢出必须提升到精确表示或报告资源错误，不能产生环绕后的证明结果。

\subsection{操作码总表}

\begin{longtable}{@{}lll@{}}
\toprule
操作码 & RPA 助记符 & 功能 \\
\midrule
\endfirsthead
\toprule
操作码 & RPA 助记符 & 功能 \\
\midrule
\endhead
\code{0x00} & \code{nop} & 无状态变化 \\
\code{0x01} & \code{ret} & 完成一个执行方向 \\
\code{0x02} & \code{drop} & 释放当前 scope 的派生事实 \\
\code{0x10} & \code{lin.le} & 线性组合得到 LE \\
\code{0x11} & \code{lin.eq} & 线性组合得到 EQ \\
\code{0x12} & \code{eq.pair} & 由相反方向不等式得到 EQ \\
\code{0x13} & \code{box.le} & 由声明变量界得到 LE \\
\code{0x14} & \code{weaken.le} & 放松一个 LE \\
\code{0x15} & \code{close.false} & 以常数矛盾关闭当前域 \\
\code{0x20} & \code{int.aff} & 组合整数事实得到 INT \\
\code{0x21} & \code{round.hi} & 整数表达式上界取整 \\
\code{0x22} & \code{round.lo} & 整数表达式下界取整 \\
\code{0x30} & \code{map.copy} & 定义输出变量的复制映射 \\
\code{0x31} & \code{map.const} & 定义输出变量的常数映射 \\
\code{0x32} & \code{map.aff} & 定义输出变量的仿射映射 \\
\code{0x40} & \code{check.row.hi} & 解除输出行上界义务 \\
\code{0x41} & \code{check.row.lo} & 解除输出行下界义务 \\
\code{0x42} & \code{check.row.eq} & 解除输出等式行义务 \\
\code{0x43} & \code{check.var.hi} & 解除输出变量上界义务 \\
\code{0x44} & \code{check.var.lo} & 解除输出变量下界义务 \\
\code{0x45} & \code{check.var.int} & 解除输出变量整数性义务 \\
\code{0x46} & \code{check.obj.eq} & 解除目标相等义务 \\
\code{0x47} & \code{check.obj.le} & 解除目标不增义务 \\
\code{0x50} & \code{split.bin} & 对显式 binary 输入变量分类 \\
\code{0x51} & \code{split.int.thr} & 对整数表达式按阈值分类 \\
\code{0x52} & \code{split.int.range} & 枚举整数表达式的有限区间 \\
\code{0x53} & \code{split.sign} & 对实仿射表达式按符号分类 \\
\code{0x54} & \code{yield} & 提交一个可达分支的汇合事实 \\
\bottomrule
\end{longtable}

\code{0x80--0xBF} 为 batch/fused 扩展保留，\code{0xF0} 为 theorem 调用保留。这些编码不属于本章的核心指令集。未声明、未支持或未知操作码必须拒绝，不能当作 \code{nop} 跳过。

\section{一般指令}

\subsection{NOP：无操作}
\label{ins:nop}

\begin{lstlisting}
nop
\end{lstlisting}

编码为 \code{opcode=0x00}、\code{flags=0}、\code{aux=0}、\code{payload_bytes=0}。NOP 不读取事实，也不生成事实、映射或义务；执行继续到下一条指令。非零保留字段或载荷以 \code{E_FORMAT} 拒绝。NOP 不能替代 root 或分支所要求的终结指令。

\subsection{RET：完成当前方向}
\label{ins:ret}

\begin{lstlisting}
ret
\end{lstlisting}

编码为 \code{opcode=0x01}、\code{flags=0}、\code{aux=0}、\code{payload_bytes=0}。RET 必须位于当前方向 root tape 的最后一条指令位置，且 scope stack 必须只包含 root。

对仍可达的 root，RET 检查：所有输出变量都有完整定义；所有要求的输出行、变量界及整数性义务均为 DONE；contract 要求的目标关系已完成。完整定义包括可用的单一表达式映射，以及经所有可达分支封存、且不再有依赖义务的分段映射。不得把封存状态误判为未映射状态。

成功后该方向成为 VERIFIED。缺失映射以 \code{E_MAP} 拒绝，未完成义务以 \code{E_OBLIGATION} 拒绝，错误 scope 或非终结位置以 \code{E_SCOPE} 或 \code{E_FORMAT} 拒绝。RET 不接受分支事实返回，也不能结束嵌套分支。root 已被证明不可达时采用机器模型章规定的空可行域结束规则。

\subsection{DROP：释放事实区间}
\label{ins:drop}

\begin{lstlisting}
drop %100 .. %499
\end{lstlisting}

编码为 \code{opcode=0x02}、\code{flags=0}、\code{aux=0}，载荷为 8 字节：

\begin{lstlisting}
first: u32;
count: u32;
\end{lstlisting}

语义作用于半开区间 \([\mathrm{first},\mathrm{first}+\mathrm{count})\) 内的派生 FactId。\code{count} 必须大于零。整个区间必须无溢出、全部已定义、全部活跃，且全部属于当前 scope。检查全部通过后一次性标记为 DEAD，随后任何事实引用都不能重新激活这些编号。DROP 不释放输入事实或分支假设。

\proposal{P-03}{为避免一个分支删除兄弟分支仍需读取的祖先事实，v0.1 草案限制 DROP 只能释放当前 scope 所拥有的派生事实。该限制是待确认的 scope 管理规则。}

空区间或算术越界以 \code{E_FORMAT} 拒绝；未定义或已死亡的事实以 \code{E_REF} 拒绝；释放祖先或其他 scope 的事实以 \code{E_SCOPE} 拒绝。DROP 改变逻辑活跃状态，但实现能否回收物理槽位取决于其事实存储结构，不能仅由该指令推出峰值内存界。

\section{线性事实指令}

\subsection{LIN.LE：非负不等式组合}
\label{ins:lin-le}

\begin{lstlisting}
%7 = lin.le $e31 {
    +1   @in.r[17].hi
    +2   %3
    -3/5 @in.r[91].eq
}
\end{lstlisting}

编码为 \code{opcode=0x10}、\code{flags=0}、\code{aux=n}。载荷长度必须为 \(4+8n\) 字节：

\begin{lstlisting}
out_expr: u32;
terms[n]: { fact: FactRef; multiplier: ScalarRef; };
\end{lstlisting}

每个输入必须为 LE 或 EQ。LE 输入的系数必须非负，EQ 输入的系数可为任意有理数。所有引用即使系数为零也必须有效，所有 scalar 也必须具有合法 canonical 编码。令输入表达式为 \(e_i\)、系数为 \(a_i\)，kernel 检查
\[
 e_{\mathrm{out}}=\sum_{i=0}^{n-1}a_i e_i.
\]
成功生成 \(\operatorname{LE}(e_{\mathrm{out}})\)。\(n=0\) 合法，但只可生成零表达式的 LE。重复输入引用按其实际出现次数参与组合，不改变推理规则。

输入 INT 以 \code{E_KIND} 拒绝；LE 的负系数、组合结果不相等以 \code{E_PROOF} 拒绝；非法引用以 \code{E_REF} 拒绝，非法 scalar 以 \code{E_SCALAR} 拒绝。

\subsection{LIN.EQ：等式线性组合}
\label{ins:lin-eq}

\begin{lstlisting}
%8 = lin.eq $e42 {
    +1 @in.r[12].eq
    -3 @in.r[19].eq
}
\end{lstlisting}

编码为 \code{opcode=0x11}、\code{flags=0}、\code{aux=n}；载荷字段与 LIN.LE 相同，长度为 \(4+8n\) 字节。所有输入必须为 EQ，所有系数可为任意精确有理数。检查
\[
 e_{\mathrm{out}}=\sum_i a_i\operatorname{expr}(r_i)
\]
后生成 \(\operatorname{EQ}(e_{\mathrm{out}})\)。空组合仅能产生 \(\operatorname{EQ}(0)\)。非 EQ 输入以 \code{E_KIND} 拒绝，结果不相等以 \code{E_PROOF} 拒绝。

\subsection{EQ.PAIR：相反不等式构成等式}
\label{ins:eq-pair}

\begin{lstlisting}
%9 = eq.pair $e17, %3, %8
\end{lstlisting}

编码为 \code{opcode=0x12}、\code{flags=0}、\code{aux=0}，载荷为 12 字节：

\begin{lstlisting}
out_expr:      u32;
positive_fact: FactRef;
negative_fact: FactRef;
\end{lstlisting}

两个输入必须为 LE，并且分别满足
\[
 \operatorname{expr}(r_+)=e_{\mathrm{out}},\qquad
 \operatorname{expr}(r_-)=-e_{\mathrm{out}}.
\]
成功生成 \(\operatorname{EQ}(e_{\mathrm{out}})\)。不接受只差非单位比例的表达式；这种情况必须先由 LIN.LE 显式调整。类型错误以 \code{E_KIND} 拒绝，任一精确比较失败以 \code{E_PROOF} 拒绝。

\subsection{BOX.LE：声明变量界检查}
\label{ins:box-le}

\begin{lstlisting}
%10 = box.le $e73
\end{lstlisting}

编码为 \code{opcode=0x13}、\code{flags=0}、\code{aux=0}，载荷仅含 \code{out_expr:u32}，长度为 4 字节。对
\[
 e=c+\sum_i a_i x_i
\]
使用当前输入模型\emph{声明}的界 \(l_i,u_i\)，计算
\[
 U=c+\sum_{a_i>0}a_i u_i+\sum_{a_i<0}a_i l_i.
\]
若每个实际使用的界均有限且 \(U\leq0\)，则生成 \(\operatorname{LE}(e)\)。canonical 表达式不含零系数，因此不存在以 \(0\cdot\infty\) 绕过检查的情况。常数表达式直接检查 \(c\leq0\)。

本指令不读取派生变量界，也不使用 branch assumption 自动收紧界。需要这些事实时应使用 LIN.LE。所需界不存在、无限或计算结果为正时以 \code{E_PROOF} 拒绝；非法变量编号或表达式编号以 \code{E_REF} 拒绝。

\subsection{WEAKEN.LE：常数放松}
\label{ins:weaken-le}

\begin{lstlisting}
%11 = weaken.le $e91, %7, 3/2
\end{lstlisting}

编码为 \code{opcode=0x14}、\code{flags=0}、\code{aux=0}，载荷为 12 字节：

\begin{lstlisting}
out_expr: u32;
fact:     FactRef;
delta:    ScalarRef;
\end{lstlisting}

输入必须为 LE，且 \(\delta\geq0\)。kernel 检查
\[
 e_{\mathrm{out}}=\operatorname{expr}(r)-\delta
\]
后生成 \(\operatorname{LE}(e_{\mathrm{out}})\)。\(\delta=0\) 合法。非 LE 输入以 \code{E_KIND} 拒绝；负 \(\delta\) 或表达式不匹配以 \code{E_PROOF} 拒绝。

\subsection{CLOSE.FALSE：关闭矛盾域}
\label{ins:close-false}

\begin{lstlisting}
close.false %27
\end{lstlisting}

编码为 \code{opcode=0x15}、\code{flags=0}、\code{aux=0}，载荷仅含 \code{fact:FactRef}，长度为 4 字节。输入必须是无变量项的常数事实，且满足以下任一条件：
\[
 \operatorname{LE}(c)\text{ 且 }c>0,
 \qquad
 \operatorname{EQ}(c)\text{ 且 }c\neq0.
\]
成功标记当前分支或 root 为 UNREACHABLE，不产生事实，不再要求当前域完成映射、义务或 YIELD。它必须是该显式指令块的最后一条指令。

非 LE/EQ 类型以 \code{E_KIND} 拒绝；非常数表达式、\(\operatorname{LE}(c\leq0)\) 或 \(\operatorname{EQ}(0)\) 以 \code{E_PROOF} 拒绝。\(\operatorname{INT}(1/2)\) 不是本指令接受的矛盾形式。root 关闭表示输入可行域为空，该方向的全称映射义务因而成立；另一个方向仍须独立验证。

\section{整数事实指令}

\subsection{INT.AFF：整数仿射组合}
\label{ins:int-aff}

\begin{lstlisting}
%12 = int.aff $e104, 7 {
    +3 %3
    -2 %7
}
\end{lstlisting}

编码为 \code{opcode=0x20}、\code{flags=0}、\code{aux=n}，载荷长度为 \(8+8n\) 字节：

\begin{lstlisting}
out_expr:       u32;
integer_offset: ScalarRef;
terms[n]: { fact: FactRef; multiplier: ScalarRef; };
\end{lstlisting}

offset 必须为整数。INT 输入的 multiplier 必须为整数；EQ 输入的 multiplier 可为任意有理数。LE 输入不允许。kernel 验证
\[
\begin{aligned}
 e_{\mathrm{out}}={}&k+
 \sum_{\operatorname{kind}(r_i)=\mathrm{INT}}n_i\operatorname{expr}(r_i)\\
 &+\sum_{\operatorname{kind}(r_j)=\mathrm{EQ}}\mu_j\operatorname{expr}(r_j),
\end{aligned}
\]
其中 \(k,n_i\in\mathbb Z\)、\(\mu_j\in\mathbb Q\)。成功生成 \(\operatorname{INT}(e_{\mathrm{out}})\)。空组合可生成任意整数常数的 INT。

\proposal{P-04}{原始草案只允许 INT 输入。本手册提出允许 EQ 项参与 INT.AFF：EQ 项恒等于零，可在不引入新操作码的情况下把整数性沿已证明等式传递，例如由 INT(z) 与 EQ(x-z) 推出 INT(x)。}

LE 输入以 \code{E_KIND} 拒绝；非整数 offset、INT 项的非整数系数或结果不相等以 \code{E_PROOF} 拒绝。输入表达式常数项本身可以是有理数，只要输入 INT 事实已经成立；kernel 不得把“所有表达式系数必须是整数”作为此指令的额外限制。

\subsection{ROUND.HI：整数上界取整}
\label{ins:round-hi}

\begin{lstlisting}
%13 = round.hi $e_x_minus_7, %4, %7
\end{lstlisting}

编码为 \code{opcode=0x21}、\code{flags=0}、\code{aux=0}，载荷为 12 字节：

\begin{lstlisting}
out_expr:  u32;
bound_fact: FactRef;
int_fact:   FactRef;
\end{lstlisting}

令 \(g=\operatorname{expr}(r_{\mathrm{int}})\)、\(b=\operatorname{expr}(r_{\mathrm{bound}})\)。输入类型必须分别为 INT 和 LE，并且 \(g-b\) 必须是常数 \(\beta\)。kernel 检查
\[
 b=g-\beta,\qquad e_{\mathrm{out}}=g-\lfloor\beta\rfloor.
\]
成功生成 \(\operatorname{LE}(e_{\mathrm{out}})\)。\(g\) 可为任意已证明整数的仿射表达式，不局限于单变量；其常数项也参与相减。取整是数学 floor，对负数同样向负无穷方向取整。

类型不匹配以 \code{E_KIND} 拒绝；\(g-b\) 含变量项或输出不匹配以 \code{E_PROOF} 拒绝。例：\(g=x\)、\(b=x-37/5\) 得到 \(x-7\leq0\)。

\subsection{ROUND.LO：整数下界取整}
\label{ins:round-lo}

\begin{lstlisting}
%14 = round.lo $e_neg_x_plus_8, %5, %7
\end{lstlisting}

编码为 \code{opcode=0x22}、\code{flags=0}、\code{aux=0}；载荷字段与 ROUND.HI 相同，长度为 12 字节。输入 \(r_{\mathrm{bound}}\) 必须为 LE，\(r_{\mathrm{int}}\) 必须为 INT。令它们的表达式为 \(b,g\)，则 \(b+g\) 必须是常数 \(\beta\)，且
\[
 b=-g+\beta,\qquad e_{\mathrm{out}}=-g+\lceil\beta\rceil.
\]
成功生成 \(\operatorname{LE}(e_{\mathrm{out}})\)。数学 ceil 对负数向正无穷方向取整。例：\(-x+37/5\leq0\) 与 \(\operatorname{INT}(x)\) 得到 \(-x+8\leq0\)。类型不匹配以 \code{E_KIND} 拒绝；和不为常数或输出不匹配以 \code{E_PROOF} 拒绝。

\section{映射指令}

\subsection{映射指令的公共条件}

本节指令只定义输出变量的取值表达式，不自动生成 LE、EQ、INT，也不解除任何输出义务。载荷中的输入变量与表达式都属于当前方向的输入坐标系。

tape 中的映射指令只能赋值给 map section 中声明为 DEFER、且在当前 scope 状态中为 UNASSIGNED 的输出变量。父 scope 已有 KNOWN 映射的变量不能被子分支重定义。各兄弟分支从同一父状态开始，因此可以分别赋值给同一个尚未定义的 DEFER 变量。重复赋值、非 DEFER 赋值或重写封存的分段映射均以 \code{E_MAP} 拒绝。

分支映射在汇合时的相等性、封存与义务依赖规则见机器模型章；任一 CHECK 需要 KNOWN 映射才能实际完成代换。

\subsection{MAP.COPY：复制输入变量}
\label{ins:map-copy}

\begin{lstlisting}
map.copy @out.v[17], @in.v[9]
\end{lstlisting}

编码为 \code{opcode=0x30}、\code{flags=0}、\code{aux=0}，载荷为 8 字节：

\begin{lstlisting}
out_var: u32;
in_var:  u32;
\end{lstlisting}

成功将 \(E_{17}(x)=x_9\) 存为 KNOWN。两个索引必须分别落在输出、输入变量范围内，否则以 \code{E_REF} 拒绝。输入变量的 domain 不自动转移到输出变量；输出 bounds 和 integrality 仍必须分别检查。

\subsection{MAP.CONST：常数映射}
\label{ins:map-const}

\begin{lstlisting}
map.const @out.v[18], 0
\end{lstlisting}

编码为 \code{opcode=0x31}、\code{flags=0}、\code{aux=0}，载荷为 8 字节：

\begin{lstlisting}
out_var: u32;
value:   ScalarRef;
\end{lstlisting}

成功将 \(E_{18}(x)=\mathrm{value}\) 存为 KNOWN。value 可以是任意合法有理数；是否符合输出变量 bounds 和 domain 由 CHECK 决定。输出索引越界以 \code{E_REF} 拒绝，非法 scalar 编码以 \code{E_SCALAR} 拒绝。

\subsection{MAP.AFF：仿射映射}
\label{ins:map-aff}

\begin{lstlisting}
map.aff @out.v[19], $e37
\end{lstlisting}

编码为 \code{opcode=0x32}、\code{flags=0}、\code{aux=0}，载荷为 8 字节：

\begin{lstlisting}
out_var: u32;
expr:    u32;
\end{lstlisting}

成功将输出变量的 KNOWN 映射设为 \(E_{19}(x)=e_{37}(x)\)。输出索引或表达式编号无效以 \code{E_REF} 拒绝。表达式不得引用输出变量，不能通过 MAP.AFF 构造递归映射。

\section{模型义务检查指令}

\subsection{CHECK 的公共语义与 EXACT 模式}

CHECK 在执行时由 kernel 从输出模型和当前映射构造待证表达式 \(P(x)\)。它不接受证书自行声明目标表达式。参与代换的每个输出变量都必须为 KNOWN；缺失或 SEALED\_PIECEWISE 映射以 \code{E_MAP} 拒绝。已封存映射对应的义务必须已经完成，不得用于新的代换。

所有 CHECK 的 \code{aux=0}。普通模式 \code{flags=0}，模型对象检查载荷为

\begin{lstlisting}
out_index: u32;
fact:      FactRef;
\end{lstlisting}

共 8 字节；目标检查仅有 \code{fact:FactRef}，共 4 字节。普通模式要求事实类型与该指令一致，且 \(\operatorname{expr}(r)=P\)。成功将对应 PENDING 义务置为 DONE，不产生可引用事实。对已经 DONE 的义务再次执行 CHECK 必须以 \code{E_OBLIGATION} 拒绝；分支继承的 DONE 义务同样不能重复检查。

EXACT 模式仅置 \code{flags} 的 bit 0，其他位必须为零。模型对象检查载荷仅有 \code{out_index:u32}，共 4 字节；目标检查没有载荷。EXACT 完全由 kernel 执行以下确定性检查：

\begin{itemize}
\item EQ 检查要求 \(P\) 为零表达式。
\item LE 检查要求 \(P\) 不含变量项且其常数 \(c\leq0\)。
\item INT 检查要求 \(P\) 的常数项为整数，并且每个非零系数为整数、对应输入变量在模型中显式声明为 integer 或 binary。
\end{itemize}

\proposal{P-05}{原始草案使用 exact 汇编 pseudo-op，但未指定二进制编码。本手册提出 CHECK 的 bit 0 作为 EXACT 模式。该模式不信任 assembler 的判断，也不生成隐式 FactId；所有表达式由 kernel 重新代换并检查。}

EXACT 不进行 BOX.LE、等式消元、推理搜索或利用派生整数事实。CHECK 对一个不存在的模型侧或不属于当前 contract 的义务以 \code{E_OBLIGATION} 或 \code{E_CONTRACT} 拒绝；事实类型错误以 \code{E_KIND} 拒绝；表达式比较或 EXACT 检查失败以 \code{E_PROOF} 拒绝。成功标记具有当前分支的可行域条件，在汇合时对所有可达分支取交集。

\subsection{CHECK.ROW.HI：输出行上界}
\label{ins:check-row-hi}

\begin{lstlisting}
check.row.hi @out.r[17], %81
check.row.hi @out.r[17], exact
\end{lstlisting}

操作码为 \code{0x40}，采用本节公共 CHECK 编码。输出行 \(r\) 必须有有限上界 \(u_r\)，且该行侧必须存在于义务集合。令行系数为 \(A_{rj}\)，构造
\[
 P=\sum_j A_{rj}E_j(x)-u_r.
\]
普通模式要求 LE(P)，EXACT 模式要求 P 是非正的常数。成功解除该行的 upper-side 义务。不存在的上界、不被要求的侧或越界行号必须拒绝。规范化等式行使用 CHECK.ROW.EQ，不把等式义务隐式拆成两个独立侧。

\subsection{CHECK.ROW.LO：输出行下界}
\label{ins:check-row-lo}

\begin{lstlisting}
check.row.lo @out.r[17], %82
check.row.lo @out.r[17], exact
\end{lstlisting}

操作码为 \code{0x41}，采用公共 CHECK 编码。输出行必须有有限下界 \(l_r\)，构造
\[
 P=-\sum_j A_{rj}E_j(x)+l_r.
\]
普通模式要求 LE(P)，EXACT 模式检查常数非正。成功解除 lower-side 义务。下界不存在、行号越界或非此侧义务按公共拒绝规则处理。行常数若存在于 canonical 模型中，必须在正规化行数据里精确计入，不得重复加减。

\subsection{CHECK.ROW.EQ：输出等式行}
\label{ins:check-row-eq}

\begin{lstlisting}
check.row.eq @out.r[17], %91
check.row.eq @out.r[17], exact
\end{lstlisting}

操作码为 \code{0x42}，采用公共 CHECK 编码。输出行必须为显式 canonical equality \(\sum_jA_{rj}y_j=b_r\)。构造
\[
 P=\sum_jA_{rj}E_j(x)-b_r.
\]
普通模式要求 EQ(P)，EXACT 模式要求 \(P=0\)。成功解除该行的 equality 义务。即使两个 LE 事实分别证明了两侧，也必须先使用 EQ.PAIR 得到 EQ 后才能作为普通模式输入。

\subsection{CHECK.VAR.HI：输出变量上界}
\label{ins:check-var-hi}

\begin{lstlisting}
check.var.hi @out.v[17], %22
check.var.hi @out.v[17], exact
\end{lstlisting}

操作码为 \code{0x43}，采用公共 CHECK 编码。输出变量 \(j\) 必须有有限上界 \(u_j\)，构造 \(P=E_j(x)-u_j\)。普通模式要求 LE(P)，EXACT 模式要求非正的常数。成功解除变量 upper-bound 义务。binary 变量的 canonical bounds 也通过该指令或普通证明检查，不由 MAP 自动解除。

\subsection{CHECK.VAR.LO：输出变量下界}
\label{ins:check-var-lo}

\begin{lstlisting}
check.var.lo @out.v[17], %21
check.var.lo @out.v[17], exact
\end{lstlisting}

操作码为 \code{0x44}，采用公共 CHECK 编码。输出变量 \(j\) 必须有有限下界 \(l_j\)，构造 \(P=-E_j(x)+l_j\)。普通模式要求 LE(P)，EXACT 模式要求非正的常数。成功解除变量 lower-bound 义务。

\subsection{CHECK.VAR.INT：输出变量整数性}
\label{ins:check-var-int}

\begin{lstlisting}
check.var.int @out.v[17], %23
check.var.int @out.v[17], exact
\end{lstlisting}

操作码为 \code{0x45}，采用公共 CHECK 编码。输出变量必须显式声明为 integer 或 binary，构造 \(P=E_j(x)\)。普通模式要求 INT(P)，EXACT 模式执行本节规定的声明整数变量线性组合检查。成功解除该变量的 integrality 义务。

EXACT 能接受 \(3x_0-2x_1+7\)（\(x_0,x_1\) 均显式整数），但不能仅凭已有 EQ 或 INT 事实接受一个连续变量 \(x_2\)。后者应先生成 INT 事实并使用普通模式。对连续输出变量调用本指令没有对应义务，必须拒绝。

\subsection{CHECK.OBJ.EQ：目标相等}
\label{ins:check-obj-eq}

\begin{lstlisting}
check.obj.eq %101
check.obj.eq exact
\end{lstlisting}

操作码为 \code{0x46}，普通载荷为 \code{fact:FactRef}（4 字节），EXACT 载荷为 0 字节，\code{aux=0}。对当前方向已经统一为最小化的目标，构造
\[
 D(x)=f_{\mathrm{out}}(E(x))-f_{\mathrm{in}}(x).
\]
两个目标的所有常数 offset 都必须包括在 D 内。普通模式要求 EQ(D)，EXACT 要求 \(D=0\)。当当前方向的目标 relation 为 EQ 时，成功完成目标相等义务；当 relation 为 LE 时，它作为更强的证明完成目标不增义务。此规则不改变证书声明或调用者所要求的 contract。relation 为 NONE 时没有可解除的目标义务，本指令必须以 \code{E_OBLIGATION} 拒绝。

缺失必要映射以 \code{E_MAP} 拒绝，不存在或已完成的目标义务以 \code{E_OBLIGATION} 拒绝。证书所声明的 contract 与调用者请求不一致属于 \code{E_CONTRACT}，应在执行之前拒绝。遗漏常数 offset 的事实会因表达式不匹配而以 \code{E_PROOF} 拒绝。

\subsection{CHECK.OBJ.LE：目标不增}
\label{ins:check-obj-le}

\begin{lstlisting}
check.obj.le %101
check.obj.le exact
\end{lstlisting}

操作码为 \code{0x47}，编码与 CHECK.OBJ.EQ 相同。构造相同的 \(D=f_{\mathrm{out}}\circ E-f_{\mathrm{in}}\)。普通模式要求 LE(D)，EXACT 要求 D 为非正的常数。当前方向必须要求 LE relation，成功完成目标不增义务；relation 为 NONE 或 EQ 时以 \code{E_OBLIGATION} 拒绝。

CHECK.OBJ.LE 不能解除目标 equality 要求，即使证书在某处还证明了相反不等式；应显式生成 EQ 并使用 CHECK.OBJ.EQ。最大化模型必须先按照 contract 的目标规范化规则转换方向，不能令同一个操作码在运行时隐式反转不等号。

\section{结构化分类与汇合指令}

\subsection{SPLIT 公共格式}

SPLIT 没有跳转目标。其载荷内包含完整、有界、顺序排列的子指令块；kernel 依次检查每个块，兄弟块共享同一父状态快照。任何指令都不能跨出所属 byte range。

\code{flags=0}，\code{aux=q} 是每个可达分支必须提交的 YIELD 数量，同时也是汇合后产生的事实数量。允许 \(q=0\)：此时分支只用于映射和义务检查。每个 \code{branch_bytes} 必须大于零且为 8 的倍数，包括该块内每条指令自己的对齐填充。每个子块必须以 YIELD 或 CLOSE.FALSE 结束，RET 不可出现在子块内。终结指令之后不得在同一子块内放置额外指令。

分支假设通过 \code{FactRef} tag 7 暴露，\code{@case[0]} 指当前分支假设；更大的 depth offset 指可见祖先分类假设。它们不是派生 FactId，不需要 LOAD，也不能被 DROP。每个 split 进入时都建立独立子 scope。

\subsection{SPLIT 公共汇合语义}

每个可达分支的 YIELD 必须具有相同 arity q，并且每个位置的 FactKind 与精确表达式相同。所有分支验证完成后，在父 scope 中为每个 yield 位置生成一个新事实。

\begin{itemize}
\item 没有可达分支：父域成为 UNREACHABLE，不提交 q 个可用事实。
\item 只有一个可达分支：它的 q 个命题成为父 scope 的汇合事实。
\item 至少两个可达分支：逐位置比较命题后生成 q 个父事实。
\end{itemize}

某个分支 CLOSE.FALSE 时不需要 YIELD，也不参与义务交集或映射比较。所有分支局部派生事实在汇合后失去可见性。已完成义务按可达分支取交集；不能把不同分支分别完成的两个义务合并为全局完成。

映射只在可达分支中比较。若一个变量在部分可达分支已定义、在另一些可达分支未定义，必须立即以 \code{E_MAP} 拒绝；所有可达分支都未定义时保留 UNASSIGNED。相同的已定义映射可保留为 KNOWN；不同映射只有在每支都定义完整、且所有依赖它的义务均在每支完成时才能封存为 SEALED\_PIECEWISE，否则以 \code{E_MAP} 或 \code{E_OBLIGATION} 拒绝。封存允许 RET，但禁止后续代换。

若嵌套 split 已证明整个当前域不可达，后续字节只接受结构、编码及索引格式检查，不进行依赖于活跃事实的语义执行，也不制造可引用事实。仍须验证该块具有合法结构终结；不可达状态不能成为跳过未知操作码或越界数据的手段。隐式 FactId 编号不因该语义跳过而改变。

\subsection{SPLIT.BIN：显式二元变量分类}
\label{ins:split-bin}

\begin{lstlisting}
(%31) = split.bin @in.v[7] {
case 0:
    %27 = lin.eq $zero { }
    yield %27
case 1:
    %30 = lin.eq $zero { }
    yield %30
}
\end{lstlisting}

本例的表达式别名 \code{zero} 表示零表达式，因此两支均返回 EQ(0)。一般情况下，两个分支的 YIELD 必须证明相同命题，不能仅因引用写法相同而通过汇合。

操作码为 \code{0x50}、\code{aux=q}，载荷固定头 16 字节：

\begin{lstlisting}
var_id:        u32;
branch0_bytes: u32;
branch1_bytes: u32;
reserved:      u32;
branch0[branch0_bytes];
branch1[branch1_bytes];
\end{lstlisting}

\code{payload_bytes} 必须等于 \(16+B_0+B_1\)。输入变量必须显式声明为 binary，不能只凭派生 INT 与界推断其 binary domain。case 0 假设为 \(\operatorname{EQ}(x_v)\)，case 1 为 \(\operatorname{EQ}(x_v-1)\)。这两支覆盖输入声明的二元域。

变量越界以 \code{E_REF} 拒绝，非 binary domain 以 \code{E_KIND} 拒绝，子块长度或保留字段错误以 \code{E_FORMAT} 拒绝。汇合类型、表达式、arity 不一致以 \code{E_BRANCH} 拒绝。输入模型对 binary 的 canonical domain 定义必须包含取值限于 \(\{0,1\}\) 的语义，不能仅把 binary 作为无约束整数标签。

\subsection{SPLIT.INT.THR：整数阈值分类}
\label{ins:split-int-thr}

\begin{lstlisting}
(%40) = split.int.thr %7, 13 {
case le:
    yield %38
case ge:
    yield %39
}
\end{lstlisting}

操作码为 \code{0x51}、\code{aux=q}，载荷固定头 16 字节：

\begin{lstlisting}
int_fact:      FactRef;
threshold:     ScalarRef;
branch0_bytes: u32;
branch1_bytes: u32;
branch0[branch0_bytes];
branch1[branch1_bytes];
\end{lstlisting}

\code{payload_bytes} 为 \(16+B_0+B_1\)。输入必须为 \(\operatorname{INT}(g)\)，threshold 必须为整数 k。第一支 \code{case le} 假设为 \(\operatorname{LE}(g-k)\)，第二支 \code{case ge} 假设为 \(\operatorname{LE}(-g+k+1)\)。它们分别代表 \(g\leq k\) 和 \(g\geq k+1\)，无需上下界即可覆盖全部整数。

非 INT 引用以 \code{E_KIND} 拒绝，非整数 threshold 以 \code{E_PROOF} 拒绝。\(k+1\) 必须精确计算，即使 k 无法放入宿主固定宽度整数。其他结构与汇合错误按 SPLIT 公共规则处理。

\subsection{SPLIT.INT.RANGE：整数有限区间枚举}
\label{ins:split-int-range}

\begin{lstlisting}
(%50) = split.int.range %int_x, %lo_x, %hi_x, 3, 7 {
case 3: yield %43
case 4: yield %44
case 5: yield %45
case 6: yield %46
case 7: yield %47
}
\end{lstlisting}

操作码为 \code{0x52}、\code{aux=q}，载荷固定头 24 字节：

\begin{lstlisting}
int_fact:    FactRef;
lo_fact:     FactRef;
hi_fact:     FactRef;
lower:       ScalarRef;
upper:       ScalarRef;
branch_count: u32;
branch_bytes[branch_count]: u32;
table_padding: zero bytes to 8-byte alignment;
branches: concatenated complete instruction blocks;
\end{lstlisting}

令 \(N=\mathrm{branch\_count}\)，表区结束地址向上对齐至 8 字节；因固定头为 24 字节，表填充长度为 \((8-(24+4N)\bmod8)\bmod8\)。载荷总长度必须为
\[
 \operatorname{align}_8(24+4N)+\sum_{i=0}^{N-1}B_i.
\]
upper 和 lower 必须为整数 \(U,L\)，并且 \(L\leq U\)、\(N=U-L+1\)。N 必须能以 u32 表示，并满足文件和执行资源限制。

输入必须精确具有以下类型与表达式：
\[
 \operatorname{INT}(g),\qquad
 \operatorname{LE}(-g+L),\qquad
 \operatorname{LE}(g-U).
\]
不接受更宽松或缩放后的界，assembler 可以先通过核心指令构造要求的精确事实。第 i 个子块对应 \(k=L+i\)，分支假设为 \(\operatorname{EQ}(g-k)\)。没有可省略的 case，也没有默认分支。允许 \(L=U\)，此时只有一个子块。

类型错误以 \code{E_KIND} 拒绝；非整数端点、逆序区间、界表达式不匹配以 \code{E_PROOF} 拒绝；branch count、表 padding、累加长度或块边界不一致以 \code{E_FORMAT} 拒绝。资源不足不能降级为少枚举若干分支。

\subsection{SPLIT.SIGN：实表达式符号分类}
\label{ins:split-sign}

\begin{lstlisting}
() = split.sign $e {
case le:
    map.aff @out.v[7], $neg_e
    yield
case ge:
    map.aff @out.v[7], $e
    yield
}
\end{lstlisting}

本例只展示不同分支映射的语法；所有依赖输出变量 7 的义务必须在各支 YIELD 之前完成。

操作码为 \code{0x53}、\code{aux=q}，载荷固定头 16 字节：

\begin{lstlisting}
expr:          u32;
branch0_bytes: u32;
branch1_bytes: u32;
reserved:      u32;
branch0[branch0_bytes];
branch1[branch1_bytes];
\end{lstlisting}

\code{payload_bytes} 为 \(16+B_0+B_1\)。令表达式为 g，\code{case le} 假设为 \(\operatorname{LE}(g)\)，\code{case ge} 假设为 \(\operatorname{LE}(-g)\)。两支覆盖全部实数；\(g=0\) 同时属于两支，允许分段映射在该交集上给出不同但分别合法的输出解。

不要求 g 为整数或预先得到其他事实。表达式编号越界以 \code{E_REF} 拒绝；保留字段、子块格式错误按公共格式规则拒绝。两支重叠不允许以单支验证替代另一支。

\subsection{YIELD：分支事实返回}
\label{ins:yield}

\begin{lstlisting}
yield %17, %19
yield
\end{lstlisting}

编码为 \code{opcode=0x54}、\code{flags=0}、\code{aux=n}，载荷为 \(4n\) 字节：

\begin{lstlisting}
refs[n]: FactRef;
\end{lstlisting}

YIELD 必须是一个子块的最后一条指令，不允许出现在 root。n 必须等于所属 SPLIT 的 q。每个引用必须有效、活跃且在当前分支可见；允许返回祖先事实、输入事实或当前 case 事实，只要最终汇合满足命题一致条件。静态输入事实的表达式可直接由不可变输入模型表示，汇合不要求在表达式池中查找或插入一个相同 ExprId。

YIELD 不在子 scope 中生成新的事实；它记录有序命题列表供父 split 比较，并结束该分支。空 YIELD 用于 q=0。父 split 完成之前，这些命题不能被另一分支引用。关闭为 UNREACHABLE 的分支不需要执行 YIELD。

非分支位置或开放 scope 错误以 \code{E_SCOPE} 拒绝；不位于块末尾或载荷长度错误以 \code{E_FORMAT} 拒绝；arity 不匹配以 \code{E_BRANCH} 拒绝；非法、dead 或超出 scope 的引用以 \code{E_REF} 拒绝。

\section{逐指令载荷尺寸速查}

下表中的尺寸不包括 8 字节指令头和指令尾部的对齐填充。\(n\) 表示 \code{aux}，\(B_i\) 表示完整子指令块字节数。所有尺寸公式均必须经过无溢出的边界检查。

\begin{longtable}{@{}lll@{}}
\toprule
指令 & \code{aux} & \code{payload_bytes} \\
\midrule
\endfirsthead
\toprule
指令 & \code{aux} & \code{payload_bytes} \\
\midrule
\endhead
NOP、RET & 0 & 0 \\
DROP & 0 & 8 \\
LIN.LE、LIN.EQ & term count & \(4+8n\) \\
EQ.PAIR & 0 & 12 \\
BOX.LE & 0 & 4 \\
WEAKEN.LE & 0 & 12 \\
CLOSE.FALSE & 0 & 4 \\
INT.AFF & term count & \(8+8n\) \\
ROUND.HI、ROUND.LO & 0 & 12 \\
MAP.COPY、MAP.CONST、MAP.AFF & 0 & 8 \\
CHECK.ROW / CHECK.VAR，普通 & 0 & 8 \\
CHECK.ROW / CHECK.VAR，EXACT & 0 & 4 \\
CHECK.OBJ，普通 & 0 & 4 \\
CHECK.OBJ，EXACT & 0 & 0 \\
SPLIT.BIN、SPLIT.INT.THR、SPLIT.SIGN & yield arity & \(16+B_0+B_1\) \\
SPLIT.INT.RANGE & yield arity & \(\operatorname{align}_8(24+4N)+\sum_iB_i\) \\
YIELD & reference count & \(4n\) \\
\bottomrule
\end{longtable}

每个 opcode 的 arity、类型与命题检查仍以对应参考条目为准。载荷尺寸正确只是进入执行语义检查的必要条件，并不构成证明成功。

\chapter{MCM 与 RPX 二进制格式}
\label{ch:format}

当前只接受 MCM 1 与 RPX 11。格式完整定义见 \code{docs/FORMAT.md}；
本章给出与执行规则对应的结构摘要。所有 framing 整数为无符号小端整数。

\section{公共记录}
16 字节头依次为四字节 magic（MCM 或 RPX 加零字节）、u16 版本、u16 零保留位、
u64 完整文件长度。不允许尾随字节。tag 为一个 u8 后接三个零字节；vector 为 u32 数量后接各项。
名称为 u32 字节数及 UTF-8 字节，模型层再执行 ASCII、长度与唯一性检查。

有理数依次为符号 tag（0 非负、1 负）、u32 分子字节数、u32 分母字节数，
然后是无符号小端幅值。分母须正、分子分母互素，无冗余高零字节。
零使用正号、空分子幅值和分母一。可选有理数为 u8 0，或 u8 1 后接有理数。
仿射表达式为有理常数和 $(\mathrm{u32},q)$ 项向量；坐标严格递增且系数非零。

MCM payload 依次为模型名、变量向量、行向量、完整仿射目标。
变量为名称、u8 整数标志、可选上下界；行为名称、表达式、可选上下侧。

\section{唯一证书布局}
RPX 11 头后直接是 forward Block，再是 reverse Block，没有外围数组计数或可选标志。
Block 为 Step 向量后接一个 End。Derivation 为 Step 向量后接 DerivationEnd。
类型 tag 为 Le=0、Eq=1、Int=2。Ref 的 tag 依第\ref{sec:fact-numbering}节顺序为 0--8；
除 Exact 无载荷外，其余均后接 u32 索引。可选 Ref 为 u8 0 或 u8 1 后接完整 Ref。

\section{操作载荷}
\begin{longtable}{lp{0.73\textwidth}}
\toprule Step & tag 后的字段顺序\\\midrule\endhead
Lin & u32 id，Kind，有理常数，$(\mathrm{Ref},q)$ 向量。\\
Round & u32 id，Le Ref，Int Ref。\\
Drop & u32 id 向量。\\
EqPair & u32 id，positive Ref，negative Ref。\\
Cases & u32 id，Kind，结论 Expr，split Expr，可选 Int Ref，左右 Derivation。\\
Plotkin & u32 id，u32 bits，u32 distance，$(\mathrm{u32},\mathrm{u64})$ 向量，clique Ref 向量。\\
\bottomrule
\end{longtable}
Step tag 为第\ref{ch:instructions}章的 0--5。DerivationEnd 为 0 Witness(Ref) 或 1 Absurd(Ref)。
End 0 为 Return（Expr 向量、Ref 向量）；1 为 Split（Expr、可选 Int Ref、左右 Block）；
2 为 Absurd（Ref）；3--7 分别为 Flow、Interval、Menu、Signed、Subflow 的 witness 后接 body Block。
旧端点 tag 8、9 拒绝。

Flow/Subflow witness 为 u32 节点数、$(\mathrm{u32},\mathrm{u32})$ 弧向量及组向量。
各组为 u32 源、$(\mathrm{u32},q)$ 需求向量、流 Expr 向量、非负 Ref 向量、平衡 Ref 向量。
Interval 为 witness 向量，每项含 lower Expr 向量、upper Expr 向量、checks Ref 向量。
Menu 为 witness 向量，每项含 u32 最大索引、activation Expr、index Expr、checks Ref 向量。
Signed 为 witness 向量，每项含 capacity 有理数、flow Expr、activation Expr、checks Ref 向量。
具体前提顺序、生成变量和行见第\ref{ch:instructions}章。

\section{结构限制}
wire 文件最多 1 GiB；每个 vector 最多 $10^7$ 项，全部计数累计最多 $2\cdot10^7$ 项；
嵌套深度最多 64；有理幅值每个最多 1700 字节；wire 名称最多 1 MiB。
计数还须与剩余字节相容，不能仅凭数量盲目分配。编码器施加同样结构限制。
模型、文本导入和路径读取还有各自更窄的限制，见第\ref{ch:loader}章。

\chapter{Loader、接口与可信边界}\label{ch:loader}

一次验证需要调用者选定的完整 source MCM、完整 target MCM、RPX certificate
和独立 Request。当前 RPX 不携带 source/target hash；验证直接读取、检查和使用这两个实际模型。
调用者有责任把正确文件或字节绑定到这两个参数，并在调用之后按 Request 的含义解释结果。
不可信 certificate 只能提交 proof，不能自行改写调用者请求。

\section{入口、处理顺序与判定}

推荐集成入口为 crate \code{kernel_core} 的 \code{api} 模块；
\code{main.rs} 只解析参数、打印 façade 返回的 JSON 和设置 exit code。
文件入口先以受限 owned read 获得完整字节，之后依次执行：
\begin{enumerate}
\item 对 source 与 target 调用 \code{codec::decode_model}，包括全部模型结构验证。
\item 调用 \code{codec::decode_certificate} 完整解析 RPX；任何嵌套分支均不能略过。
\item 调用 \code{rpa::parse_request} 解析独立请求。
\item 核对 certificate 与 Request 的全部 claims、point claims、bijection claims
及 improvement claims 在数量、顺序与每个字段上精确相等。
\item 用实际模型验证各方向及其 complete Returns；随后验证 bijection injectivity、
point witnesses，最后验证 improvements。所有阶段共享同一次工作预算。
\item 只有全部请求通过后返回 \code{VERIFIED}。解析、导入或装配成功不具有该含义。
\end{enumerate}
模型导入可以事先完成。MPS/LP 不在 \code{verify} 内重复解析，也没有 solver 回调。

\begin{longtable}{p{32mm}r p{94mm}}
\toprule 状态 & exit & 含义 \\\midrule\endfirsthead
\toprule 状态 & exit & 含义 \\\midrule\endhead
\code{VERIFIED} & 0 & 独立 Request 的全部数学合同已由此 kernel 接受。\\
\code{CANONICALIZED} & 0 & MPS/LP 解析、模型检查与 MCM 编码完成。\\
\code{ASSEMBLED} & 0 & strict RPA 解析与 RPX 编码完成。\\
\code{DISASSEMBLED} & 0 & RPX 解码并输出 RPA。\\
\code{INSPECTED} & 0 & MCM 解码并报告规模。\\
\code{EXPORTED} & 0 & MCM 解码并导出规范模型 JSON。\\
\code{INVALID} & 2 & 格式、模型、请求匹配、证明或数学义务不满足要求，亦包括不支持的版本/方言。\\
\code{RESOURCE_LIMIT} & 3 & 已公布预算耗尽；不是验证成功，也不是数学反驳。\\
\code{IO_ERROR} & 4 & 文件读写失败。\\
\bottomrule
\end{longtable}
没有 exit 1 的规定状态，也没有单独的 \code{UNSUPPORTED}。
调用者必须看 status，不能把辅助命令的 exit 0 当成模型关系成立。
\code{INVALID} 也不统一表示“目标数学命题为假”：格式错误和 importer 能力限制都属于该状态。

\section{可信精确模型导入}

\subsection{共同规范化与模型身份}
MPS importer 位于 \code{model.rs}，LP importer 位于 \code{lp.rs}；
两者和 MCM codec 都是可信代码。它们用任意精度整数与精确有理数进行解析及规范化，
不经过浮点数，不调用优化器，不修改原始文件。
每个 row/目标表达式把同一变量系数精确相加，按 coordinate ID 排序并删除精确零。
零系数的变量出现仍声明变量，不能因为最后项取消而删除该变量。
名称、原始声明次序、所有变量/row、bounds、integer flags 和完整目标均保留于规范模型。

所有内部目标均最小化。maximization 使\emph{整个}仿射目标取负，常数也要取负。
缺失 bound side 是真正的无限方向；不能用大有限数代替。
互相矛盾的有限 bounds 是空模型的合法表示，不能为了得到可行输入而“修复”。
integer 的有理 bounds 不自动取整，因为 LP relaxation 的含义也须保留。
模型结构合法不证明模型可行，更不证明 source/target 等价。

\subsection{MPS 的 sections、cards 与排列}
MPS headers 必须从第一列开始，data cards 缩进。次序为 NAME，
可选 OBJSENSE/OBJNAME，ROWS，COLUMNS，可选 RHS、RANGES、BOUNDS，最后必须 ENDATA。
OBJSENSE 与 OBJNAME 各至多一次；single value 可放在 header 或下一张缩进 card。
senses 接受 MIN/MINIMIZE、MAX/MAXIMIZE，默认 MIN；headers、keywords、sense 区分大小写。
NAME 未给名称时使用 \code{MODEL}。

必须恰好有一个 N objective row，E/L/G 是 constraint rows。
OBJNAME 若存在必须指定唯一 N row。constraint 顺序为 ROWS 顺序，除去 objective row；
variable 顺序为首次 COLUMNS 出现顺序，重复 column cards 必须连续。
每个 coefficient card 含一或两组 row/value。重复 coefficient 可相加，
重复 row name、未知 row/variable、非连续 column 声明均拒绝。

本方言接受空白分隔 free MPS，以及名称不含空白的普通 fixed-column cards。
\code{COLUMNS/RHS/RANGES} 的 fixed continuation 只有前十四个字符全为空格时才能省略
column/vector name，而且此前须存在可继续的同类 card。
fixed-format 的 embedded-space 名称拒绝；不能靠对齐猜测缺失必要字段。
第一列 \code{*} 开始注释；独立 token 以 \code{\$} 开始时截断后续 inline comment，
附着在名称中的 \code{\$} 不会截断数据。因此本方言不接受以 \code{\$} 开始的数据名称。
空白行和注释行忽略；ENDATA 后非注释内容拒绝。

RHS、RANGES、BOUNDS 每种至多一个 named vector；不默选第一个向量。
同一 row 的重复 RHS 或 range card 拒绝。缺失 RHS 是精确零。
objective row 的 RHS 值 $b$ 生成 objective constant $-b$，然后才按 sense 归一化；
objective 上不能有 range。

\subsection{MPS bounds、markers 与 ranges}
连续 column 默认 $[0,+\infty)$。\code{'MARKER' 'INTORG'} 至
\code{'MARKER' 'INTEND'} 的 column 是 integer；markers 必须配对且不能嵌套。
标记 integer column 如果\emph{完全没有}显式 bound card，默认 $[0,1]$；
一旦存在任何显式 card，未定义 sides 回到一般 $[0,+\infty)$ defaults。
故单独 LI 表示 general integer，不再有隐式 upper 1；标记 integer 的 PL 也显式移除 upper。

\begin{longtable}{p{18mm}p{128mm}}
\toprule bound card & 精确含义 \\\midrule\endfirsthead
\toprule bound card & 精确含义 \\\midrule\endhead
LO / UP & 设置有限 lower / upper；可以为有理数。\\
FX & 两侧设为同一个有限有理数。\\
FR / MI / PL & 两侧 infinite / lower infinite / upper infinite；不允许 numeric field。\\
LI / UI & 设置 lower / upper 并声明 integer；该 card 的值必须是整数。\\
BV & integer 的 $[0,1]$ 与所有显式 bounds 相交，card 顺序无关；可选数值必须恰好为 1。\\
\bottomrule
\end{longtable}
同一 side 重复定义必须相同；冲突的覆盖拒绝。相反 sides 的矛盾保留。
若显式 upper 为负且没有显式 lower card，MPS 默认 lower 变为 $-\infty$；
upper=0 则仍保留默认 lower=0。显式 lower 永远优先于此默认规则。
BV 则始终和 $[0,1]$ 相交：BV+LO 2 得 $[2,1]$，BV+FR 得 $[0,1]$，
BV+UP $-1$ 得 $[0,-1]$，后者不会被负 upper 默认规则变为可行区间。

row RHS 为 $b$、range 为 $r$ 时：
\begin{longtable}{p{27mm}p{119mm}}
\toprule row 类型 & canonical finite sides \\\midrule\endfirsthead
\toprule row 类型 & canonical finite sides \\\midrule\endhead
L & $b-|r|\le a^\top x\le b$。\\
G & $b\le a^\top x\le b+|r|$。\\
E，$r\ge0$ & $b\le a^\top x\le b+r$。\\
E，$r<0$ & $b+r\le a^\top x\le b$。\\
\bottomrule
\end{longtable}
没有 range 时保留原 equality 或单侧 constraint；range=0 生成 equality。

MPS 明确拒绝多个 N row、多个 vectors、初始 blank-name vectors、未知 section/bound type、
SOS、quadratic/conic/indicator、semicontinuous/semiinteger 及其他未定义扩展。
拒绝表明本 importer 没有接受这种方言，不能据此判断模型间关系。

\section{LP v1 的自足文本规则}

\subsection{sections、名称与 coordinate order}
\code{canonical} 依据输入文件扩展名选择 parser：大小写不敏感的 \code{.lp}
调用 LP importer，其他扩展调用 MPS importer；没有内容 sniffing。
LP 输入必须 ASCII；反斜线开始到该物理行末的注释。
除 line endings、tab 外的 control characters 拒绝。
keywords 不区分大小写，每个 section header 独占一行。

必须先 objective，后 constraints，最后 End。Bounds 可选且位于 constraints 之后、
type sections 之前。Binary、General、Semi 可按任意顺序放在 constraints/Bounds 之后，
各至多一次；Bounds 不能出现在 type section 之后。
\begin{longtable}{p{33mm}p{113mm}}
\toprule section & 接受的 aliases \\\midrule\endfirsthead
\toprule section & 接受的 aliases \\\midrule\endhead
min objective & \code{Minimize, Minimum, Min}\\
max objective & \code{Maximize, Maximum, Max}\\
constraints & \code{Subject To, Such That, St, St., S.t.}\\
bounds & \code{Bounds, Bound}\\
binary & \code{Binary, Binaries, Bin}\\
general integer & \code{General, Generals, Gen}\\
empty semi marker & \code{Semi, Semis, Semi-Continuous}\\
end & \code{End}\\
\bottomrule
\end{longtable}
Semi 只允许空 section；任何成员拒绝。End 后只能有空行或注释。

名称以 ASCII letter 或 underscore 开始；后续可用字母、数字及
\code{_.$\#(),;!?@\{\}~\%\&}。名称区分大小写、最长 1,024 B、不截断。
以 e/E 为首，后面缺省或紧接数字、e/E 的 exponent-ambiguous name 拒绝。
名称不能是 section/type keyword、free、inf、infinity 或保留的未支持扩展 keyword：
sos、lazy、scenario/scenarios、indicators、pwl、multi-objectives。
quoted、square-bracket names 和数字/标点起首名称不支持；名称内括号是字面字符，
不是表达式 grouping。

LP 没有 problem-name field，模型名固定为 \code{LP}。
可选 objective label 只检查语法，不另存为模型字段。
variable 顺序按 objective、rows、bounds、type lists 的首次文本出现决定；
constraint 顺序按输入决定。constraint labels 保留且不能重复；
unnamed rows 自动命名为 \code{_lp_row_0}、\code{_lp_row_1} 等。
bounds/type-only 变量也保留。

\subsection{表达式、续行与关系}
表达式由 signed terms 组成：数值常数，或可选系数后接变量。
首项以后每项都要 \code{+} 或 \code{-}。系数与变量必须有空白，
例如 \code{2 x} 合法、\code{2x} 拒绝；operator 自身无需两侧空白，\code{x-y} 合法。
符号可以与 magnitude 分离；objective 和 named row 的 coefficient/variable tokens
可跨物理行。不能使用乘号、power、quadratic brackets、grouping 或函数调用。
零 objective 必须写 \code{0}，不能省略整个表达式。

若第一条 constraint 带 label，只有新 data line 开头的 \code{label:}
开始新 row，其他行继续当前 row。若第一条 constraint 不带 label，
每个物理 constraint line 必须包含一条完整 unnamed relation；
后续 named row 拒绝。不能由 parser 猜测未标号行是续行还是新 row。

单关系为 \code{affine <= affine}、\code{affine >= affine} 或
\code{affine = affine}。reader 减去右端，再把差的 constant 移到 bound。
例如 $x+1\le y+3$ 变为 $x-y\le2$。
range row 允许 \code{constant <= affine <= constant} 或反向的两个 \code{>=}；
方向必须一致，两 endpoints 必须精确化简成常数。
例如 $1\le x+2\le3$ 变为 $-1\le x\le1$；所有导入的 row Expr constant 为零。
named range row 可跨行；strict inequalities 拒绝。

\subsection{bounds 与 integrality}
continuous/general integer 默认 $[0,+\infty)$；显式 bound 替换指定 side。
LP 的负 upper \emph{不会}移除默认 lower zero，故只有 \code{x <= -1} 时为 $[0,-1]$。
这与前节 MPS convention 不同，文件扩展选择是模型语义的一部分。

每个 Bounds record 占一物理行，只约束一个未缩放 variable。
\begin{lstlisting}
0 <= x <= 3
3 >= x >= -2
x <= 4
-2 <= x
x >= -2
4 >= x
x = 1
1 = x
x free
-inf <= x <= +infinity
\end{lstlisting}
inf/infinity 不区分大小写，仅允许负 infinite lower 或正 infinite upper；
infinite equality、正 infinite lower、负 infinite upper 均拒绝。
同一 side 的重复 bound 必须相同；冲突 override 拒绝。
Binary 是 integer 的 $\{0,1\}$ 与 explicit bounds 的交集，
例如 binary x 加 $2\le x\le7$ 为 $[2,1]$，binary x free 为 $[0,1]$。
General 仅声明整数性，不隐含 upper=1。
type lists 为以空白分隔的名称；重复或冲突 type declarations 拒绝。

LP 拒绝未知 sections/tokens、SOS、indicators、scenarios、quadratic/nonlinear syntax、
非空 Semi、缺失必要 section/End、End 后 noncomment content 和未定义 dialect ambiguities。
有理 fraction、双侧 affine relations 是本 importer 明确提供的扩展；
vendor solver 接受范围不自动成为本 kernel 的输入语义。

\section{精确数值与各层资源边界}

MPS、LP 和 RPA 共用 \code{parse_q}。它接受带符号 integers、integer fractions、
有限 decimal（含 \code{.1}、\code{1.}）及 e/E/d/D scientific notation，
包括 \code{1/-2}；分母零拒绝。text 可以是未约分或负分母拼写，解析后规范化；
binary Q 则必须已经规范，不能自动约分。
NaN、infinity coefficients、hexadecimal floats、数字内部 whitespace 和 digit separators 不接受。
无限 bounds 仅由上述 importer grammar 表示。

\begin{longtable}{p{90mm}p{56mm}}
\toprule 层与资源 & 当前上限 \\\midrule\endfirsthead
\toprule 层与资源 & 当前上限 \\\midrule\endhead
façade/CLI 每个 owned input file & 512 MiB\\
RPA / Request text & 256 MiB；JSON nesting capacity 128\\
MPS / LP text & 512 MiB；每物理行 32,768 B\\
LP lexical tokens & 20,000,000\\
model names / variables / rows & 1,024 B / 2,000,000 / 2,000,000\\
MPS ROWS cards（含 objective） & 2,000,000\\
coefficient entries / final model terms & 20,000,000\\
decimal mantissa 或 fraction component digits & 每部分 4,096\\
decimal exponent digits / absolute exponent & 6 / 4,096\\
effective power of ten & 绝对值至多 4,096\\
scalar text bytes & $2\times4,096+32$\\
exact intermediate numerator / denominator & 各 65,536 bits\\
binary codec 文件 / 单 array / 累计 counted items & 1 GiB / 10,000,000 / 20,000,000\\
binary magnitude / name preflight & 每 magnitude 1,700 B / 1 MiB\\
binary proof depth & 64，root=1\\
默认 verification work / live facts / depth & $10^9$ / $10^6$ / 64\\
\bottomrule
\end{longtable}
限制叠加适用，不能用宽松的一层覆盖严格的一层。
例如 text parser 能构造的 65,536-bit Q 可能超过 MCM 的 1,700-byte magnitude 编码预算。
\code{verify_bytes} 不适用路径读取的 512 MiB cap，但仍适用 codec、Request 和验证预算。
theorem constructors 还有自己的维数和工作限制；生成模型也必须满足同一个 Model validation。
预算耗尽没有浮点 fallback、wraparound 或“尽力接受”的模式。
重要 allocations 使用 \code{try_reserve}；本接口不承诺第三方任意精度库的所有内部
allocation failures 或宿主强制终止都能恢复成 JSON。

work 在各方向、lift、injectivity、point refutation 与 improvements 间累积，
嵌套 scopes 不重置。live facts 限制当前可见事实；proof nesting 与 JSON parser nesting
是两种不同限制。CLI 可以设 work 和 live facts，不能设置 depth。
版本最低值与资源限制也不同：V10 本身不提高默认预算。

\section{真实 CLI 与 Rust façade}

当前命令如下；二进制名为 \code{rpx}，仓库构建命令为
\code{cargo build --release -p kernel-core --bin rpx}。
\begin{lstlisting}
rpx canonical input.mps output.mcm
rpx canonical input.lp output.mcm
rpx assemble input.rpa output.rpx
rpx disassemble input.rpx output.rpa
rpx inspect-model model.mcm
rpx export-model model.mcm output.json
rpx verify --work 1000000000 --live-facts 1000000 \
  source.mcm target.mcm proof.rpx request.json
\end{lstlisting}
verify 的两个 options 可出现在 subcommand 后任何位置；paths 仍须恰好四个，
顺序为 source、target、proof、request。数值分别解析为 u64 和 usize。
其他 command 不接受这些 flags，没有 \code{--profile}、\code{--direction}、
hash 参数或 certificate 自定 budget。未知 option、缺参数、路径数量错误返回 usage INVALID。

façade 的路径函数与 CLI 同名操作对应：
\code{canonicalize_path}、\code{assemble_paths}、\code{disassemble_paths}、
\code{inspect_model_path}、\code{export_model_paths}、\code{verify_paths}。
需要内存集成时使用 \code{api::verify_bytes(source,target,proof,request,limits)}，
前三项为 byte slices、request 为 text；不触及文件系统。
入口返回 \code{Result<serde_json::Value, api::Error>}。
更底层的 \code{verify::verify_request} 接受内存中的 Model/Certificate/Request 与 Limits，
并自行再次验证两个完整模型；
带位置的接口为 \code{verify_request_with_diagnostics}。
legacy \code{verify} 接受 Claim slice，只请求方向，不接受 points/bijections/improvements。
独立 LP contradiction 和完整 injectivity 库入口分别为
\code{verify_lp_refutation}、\code{verify_injectivity}，并有相应 diagnostics 版本。
低层 \code{verify_injectivity} 只检查给定完整 map 的单射性；它不独立建立总可行性、
与证书索引的绑定或完整双射。完整合约验收使用 \code{verify_request}。

路径读取先检查 metadata length，再用限长 read 确保文件增长也不超过 cap；
实际用于验证的是取得的 owned bytes。
这不是操作系统级多文件事务：如需 source/target/proof/request 的一致业务快照，
由调用者先获取并通过 \code{verify_bytes} 绑定。
output 在内容全部生成后创建父目录并写入指定文件；路径函数可能替换现有 output，
不承诺 atomic rename 或自动 backup。

\section{报告、诊断与模型 JSON}

每次 CLI invocation 向 stdout 打印一个 JSON value。verify 成功包括
status、claims/points/bijections counts、statistics、limits、
\code{model_decode_us}、\code{proof_decode_us}、\code{execute_us}、\code{total_us}；
路径调用另有 \code{read_us}，内存调用没有该字段。
非空 improvement 请求增加 \code{improvements} count；新增 V10 statistics 零值省略。
这些 timings 是本次调用诊断量，不能单独解释为隔离 benchmark。

failure 必含 status、message、code；已知时增加 location 与 limits。
\begin{lstlisting}
{"status":"INVALID",
 "message":"INVALID obligation mismatch: expected Eq; obligation 3",
 "code":"obligation_mismatch",
 "location":{"claim":0,"block":0,"obligation":3},
 "limits":{"work":1000000000,"live_facts":1000000,"depth":64}}
\end{lstlisting}
code 是错误词汇生成的 snake case 诊断辅助，不是冻结枚举；
稳定分类应以 status 为准。Location 全部 indices 为零基：
claim/point 是 Request 各自数组下标，block 是顶层 certificate Block 下标，
step 是最内层活动 Block/Derivation 的 Step 下标，
obligation 是当前 terminal 的实际检查次序下标。
没有相应语义的位置省略字段。improvement、fiber refutation 等不是顶层方向 Block，
不能从不存在的 claim/block 字段推断坐标。
diagnostics 只是 best-effort error attribution，不参与数学决定、统计或预算。
CLI 在有效 verify 参数解析后的失败上附加 effective limits；
库调用者可用 \code{Error::with_limits} 附加。

\code{inspect-model} 报告模型名、variables/rows、integer count 和 row/目标总 nonzeros；
\code{export-model} 通过 \code{model::to_json} 输出
name、variables、rows、objective。每个 rational 为 canonical scalar string，
bounds 缺侧为 null，terms 是有序的 \code{[variable,coefficient]} pairs。
JSON object key 顺序不具有语义，数组次序保留模型坐标。
这是 export schema，不能直接使用 Model 的普通 Serialize derive 作为同样的 scalar interchange，
也没有当前 CLI 的 JSON-model import command。

\section{信任结论与外部证据}

接受边界包含精确数值 parser/算术、MPS/LP importers、Model validation、
MCM/RPX codec、Request matching、scalar/scope checker、完整 obligation construction、
point 与 injectivity checker，以及各静态 lift/theorem 的数学 lemma 和 Rust 实现。
扩展的 old-scope premise checking 只证明其前提；生成模型合法且 child proof 通过，
仍不能独立证明联合存在性 lemma。这一项必须明确列入 TCB。

certificate generator、求解器、Python oracle、model-specific structure scripts、
外部 SHA receipts、benchmark 和 corpus replay 都不向 kernel 提供 proof facts。
独立测试、字节攻击、构造 oracle 与 source review 是实现可靠性的证据，
不是 Rust 实现的形式验证。V9/V10 的冻结来自 2026-10-14 的明确项目决定，
其基线与复核记录见 \code{audit/PROTOCOL_V9_V10_FREEZE.md}。
验证已准许请求不自动获得未请求的双向等价、最优性、可行性或 global inclusion。

本章以 \code{api.rs}、\code{main.rs}、\code{model.rs}、\code{lp.rs}、
\code{verify.rs} 为实现依据；输入方言详解对应 \code{docs/MPS_V1.md}、
\code{docs/LP_V1.md}，接口说明对应 \code{docs/API.md}。

\chapter{RPA 文本汇编与工具链}
\label{ch:assembly}

RPA（Reformulation Proof Assembly）是 RPX 的文本表示。文本层承担符号解析、精确数值解析、表达式展开和二进制编码；证书是否成立仍由 RPX verifier 判定。本章定义用于人工编写、生成器输出和反汇编的基础语法，不赋予 assembler 任何独立的证明权力。

本章中的命令行和工具名称是 v0.1 工具接口草案。它们描述计划交付的行为，不表示当前仓库已经实现这些程序。MCM 的二进制编码和完整 MPS 导入覆盖范围尚待单独冻结。

\section{文件结构}

\subsection{编码与词法约定}

RPA 文件采用 UTF-8 编码，不含 BOM。语言关键字、标识符和操作数字符使用 ASCII；注释可以使用 UTF-8。关键字区分大小写，规范输出采用本章给出的小写拼写。空格和制表符是 token 分隔符；行结束可以是 LF 或 CRLF。assembler 在词法处理时将 CRLF 视为 LF。

\code{;} 引入行注释，注释持续到该行结束。基础语法不支持块注释、宏、条件汇编、文件包含或用户定义伪指令。分号不是语句分隔符。除花括号中的列表外，一条语句必须在同一物理行结束；列表允许跨行。文件末尾可以省略最终换行。

标识符由字母或下划线起始，随后可以包含字母、数字和下划线。Fact 和 expression 标签另外允许全数字名称。这里的数字名称仍然是\emph{符号名称}，不是用户指定的二进制 ID。例如作者写的 \code{\%99} 可以被分配为第一个 derived FactId；规范反汇编则写成 \code{\%0}。同一 execution domain 内，同类标签不得重复定义。两个 domain 的标签命名空间独立。

源码中的变量、行和静态事实索引必须是非负十进制整数。前导零允许被 assembler 接受，规范输出必须删除前导零。引用索引必须满足对应模型或 FactRef 编码的取值范围；表达式不得引用输出变量。

\subsection{顶层声明}

一个基础 v0.1 文件包含版本、两个模型摘要、contract 和两个 execution domain。\code{.forward} 和 \code{.reverse} 各出现一次，并按此顺序排列。每个 domain 恰好包含一个 \code{.map} 区域和一个 \code{.code} 区域；所有 \code{.expr} 声明位于 \code{.map} 之前。

\begin{lstlisting}[caption={RPA 文件结构；尖括号表示说明性占位符}]
.rpa 0.1
.source sha256 "<64 hexadecimal digits>"
.target sha256 "<64 hexadecimal digits>"
.contract mip=opt-equivalent
.contract forward.obj=eq
.contract reverse.obj=eq

.forward
.expr $name = <affine expression>
.map
    <static mapping declarations>
.endmap
.code
    <proof instructions>
    ret
.endcode

.reverse
.expr $other = <affine expression>
.map
    <static mapping declarations>
.endmap
.code
    <proof instructions>
    ret
.endcode
\end{lstlisting}

\code{.source} 和 \code{.target} 的字符串必须各包含恰好 64 个十六进制数字。输入可以使用大小写字母，规范输出使用小写。摘要是相应 canonical model 的 SHA-256，不是原始 MPS 文本的摘要。hash 字符串不支持转义或空白。

\begin{longtable}{p{0.32\textwidth}p{0.60\textwidth}}
\toprule
声明 & 基础 v0.1 取值与约束 \\
\midrule
\endhead
\code{.rpa 0.1} & 必须为第一条非注释语句；一个文件只出现一次。\\
\code{.contract mip=opt-equivalent} & 要求双向可行映射与双向目标关系；每个方向的目标关系必须为 \code{eq} 或 \code{le}。\\
\code{.contract mip=feasibility-equivalent} & 要求双向可行映射；允许 \code{none}、\code{eq} 或 \code{le}。指定 EQ/LE 时还必须完成对应目标义务。\\
\code{.contract forward.obj=eq} & Forward 的归一化目标差要求等于零。\\
\code{.contract forward.obj=le} & Forward 的归一化目标差要求不大于零。\\
\code{.contract forward.obj=none} & 仅用于可行性等价 contract；不生成 Forward objective obligation。\\
\code{.contract reverse.obj=...} & 与 Forward 相同，但 input/output 模型方向交换。\\
\bottomrule
\end{longtable}

每条 contract 声明必须恰好出现一次，不能通过重复声明覆盖先前取值。Forward 和 Reverse 共同证明文件声明的 contract。调用者要求的最小 contract 由 verifier 的调用参数决定；文件不能自行降低调用者要求。最大化目标和最小化目标的归一化定义见机器模型与 contract 章节。

\section{精确数值字面量}

\subsection{允许的拼写}

基础语法接受整数、分数、有限十进制数和十进制科学计数法。分数字面量的分母必须为严格正整数；符号只放在分子之前。

\begin{lstlisting}
17
-3
4/7
-22/13
0.125
.125
1.
1e-6
-2.5E+3
\end{lstlisting}

这些写法均通过整数运算解析，不经过 IEEE 浮点数。十进制和科学计数法定义如下。将十进制尾数的全部数字拼接为整数 $d$，设小数点后有 $k$ 位，指数为整数 $t$，符号为 $s\in\{-1,1\}$，则：
\[
  \operatorname{value}=s\,d\,10^{t-k}.
\]
例如 \code{0.125} 解析为 $125/1000=1/8$，\code{1e-6} 解析为 $1/1000000$，\code{-2.5E+3} 解析为 $-2500$。分数通过约分得到 canonical rational。\code{-0}、\code{0/7} 和 \code{0.0} 均归一化为整数零。

\must{assembler 必须以精确整数运算实现数值字面量转换，不能先解析为浮点数再尝试恢复有理数。}

基础语法不接受 \code{NaN}、\code{Inf}、十六进制数、数字分组下划线、循环小数或隐含分母符号。\code{1/0} 是错误；\code{1/-2} 必须写为 \code{-1/2}。资源限额可以限制字面量长度和指数绝对值，但不得改变其数学值。

\subsection{规范输出}

规范反汇编只输出整数或约分后的 \code{p/q}：分母严格正，整数不带 \code{/1}，正数不带前导加号，零写作 \code{0}。因此源码中的十进制拼写、科学计数法、前导零和原始分数形式不会被保留。

\section{操作数与符号}

\begin{longtable}{p{0.29\textwidth}p{0.28\textwidth}p{0.35\textwidth}}
\toprule
操作数形式 & 类别 & 意义 \\
\midrule
\endhead
\code{\%name} & Derived FactRef & 当前 domain 已定义且可见的事实。\\
\code{@in.r[n].hi} & 静态 FactRef & input row 的有限 upper side，类型为 LE。\\
\code{@in.r[n].lo} & 静态 FactRef & input row 的有限 lower side，类型为 LE。\\
\code{@in.r[n].eq} & 静态 FactRef & input equality row，类型为 EQ。\\
\code{@in.v[n].hi} & 静态 FactRef & input variable 的有限 upper bound，类型为 LE。\\
\code{@in.v[n].lo} & 静态 FactRef & input variable 的有限 lower bound，类型为 LE。\\
\code{@in.v[n].int} & 静态 FactRef & input variable 的 integrality，类型为 INT。\\
\code{@case} & Branch FactRef & 当前最内层分支假设；等价于 \code{@case[0]}。\\
\code{@case[n]} & Branch FactRef & 距当前分支 $n$ 层的祖先分支假设。\\
\code{\$name} & ExprId 符号 & 当前 domain 的 canonical affine expression。\\
\code{@in.v[n]} & Input VarId & 表达式或 mapping 中的 input variable。\\
\code{@out.v[n]} & Output VarId & mapping 或 variable obligation 的对象。\\
\code{@out.r[n]} & Output RowId & row obligation 的对象。\\
\bottomrule
\end{longtable}

静态 FactRef 不产生指令，也不占用 derived FactId。只有模型确实声明的事实可以被引用：没有有限 upper bound 时，\code{@in.v[n].hi} 不存在；连续变量没有 \code{.int} 事实；非 equality row 没有 \code{.eq} 事实。equality row 的 side 可用性遵循 canonical model 对有限 sides 的定义。

\code{@case[n]} 中的 $n$ 是 lexical scope depth offset，不是分支编号。当不存在该层假设时引用无效。\code{@case} 的具体类型和表达式由 split 指令规定，不能由源码重声明或赋值。

\section{表达式声明与展开}

\subsection{仿射语法}

表达式语法支持精确常数、input variable、此前声明的 expression 符号、括号、加减法及标量乘法：

\begin{lstlisting}
.expr $x = @in.v[0]
.expr $sum = @in.v[0] + @in.v[2] - 1
.expr $scaled = 3/5 * $sum + 2
.expr $neg = -$scaled
.expr $mixed = 2 * (@in.v[0] - @in.v[2]) + 7
\end{lstlisting}

expression 符号通过展开得到 affine expression；符号不是运行时节点。引用必须指向同一 domain 内此前声明的 expression，因此展开关系无环。每个声明最终归一化为一个常数和严格按 VarId 排序的非零系数列表。

不允许 variable 与 variable 相乘、不允许两个一般表达式相乘，也不允许除以一般表达式。乘法的左侧必须是一个数值字面量。\code{3 * @in.v[0]} 合法；\code{@in.v[0] * 3} 不属于基础语法。减法、负号和括号可用于任意 affine expression。\code{2 @in.v[0]} 这样的隐含乘法不合法。

\subsection{指令中的 expression 操作数}

基础指令的 expression 操作数必须为 \code{\$name}。例如 \code{map.aff @out.v[0], -\$e} 不合法，必须先声明 \code{.expr \$neg = -\$e}，再使用 \code{map.aff @out.v[0], \$neg}。这一限制确保 instruction 编码直接引用 expression pool，并让反汇编语法唯一。

assembler 可以合并数值相同的 expression 或 scalar pool 条目；它必须保留程序的数学语义和 scope 行为。基础文件格式是否允许重复 pool 条目由 wire format 章节规定，文本标签不承担 pool 唯一性保证。

\section{静态 mapping 区域}

\code{.map} 描述不依赖分支的 output variable mapping。它的变量命名方向随 domain 改变：Forward 中 input 是 Source、output 是 Target；Reverse 中相反。

\begin{lstlisting}
.map
    copy.range @out.v[0], @in.v[0], 1000
    copy @out.v[1000], @in.v[7]
    const @out.v[1001], 0
    aff @out.v[1002], $replacement
    defer @out.v[1003]
    defer.range @out.v[1004], 4
.endmap
\end{lstlisting}

\begin{longtable}{p{0.18\textwidth}p{0.39\textwidth}p{0.35\textwidth}}
\toprule
语句 & 文本操作数 & section 编码与含义 \\
\midrule
\endhead
\code{copy.range} & output 起点、input 起点、count & COPY\_RANGE；连续 $count$ 对变量作 identity-copy。\\
\code{copy} & output variable、input variable & COPY\_PAIRS 的单对声明；assembler 可以合并多条声明。\\
\code{const} & output variable、number & CONST；output 等于该精确常数。\\
\code{aff} & output variable、expression & AFF；output 等于该 input affine expression。\\
\code{defer} & output variable & DEFER；该变量必须在 proof tape 的可达路径上定义。\\
\code{defer.range} & output 起点、count & 连续 DEFER 声明的文本缩写。\\
\bottomrule
\end{longtable}

\must{静态 map 必须恰好覆盖全部 output variables；每个变量恰好由 COPY、CONST、AFF 或 DEFER 覆盖一次。}

range 的 count 必须为严格正整数，索引范围的加法必须检查溢出。两个范围重叠、copy 越界或重复声明均是错误。一个静态 DEFER 标记不是已经完成的 mapping，也不能消除任何 domain obligation。

\code{map.copy}、\code{map.const} 和 \code{map.aff} 是 proof tape 指令，语法与相应静态声明相近。它们只用于当前路径尚未赋值的 DEFER output variable。对已有 mapping 重复赋值不合法。不同 sibling branches 可以各自给同一个 DEFER variable 赋值，因为两个分支分别从父状态执行。

\section{基础指令的文本形式}

下列语法穷尽基础 v0.1 的 28 个 core opcodes。所有指令的数学检查、异常条件和二进制 payload 见 ISA 章节；本节只规定文本 spelling 和操作数顺序。

\subsection{General、linear 与 integer 指令}

\begin{lstlisting}
nop
ret
drop %fact
drop %first .. %last

%out = lin.le $expr {
    +1 @in.r[0].hi
    -3/5 @in.r[1].eq
}
%out = lin.eq $expr {
    +1 @in.r[0].eq
    -2 %earlier
}
%out = eq.pair $expr, %positive, %negative
%out = box.le $expr
%out = weaken.le $expr, %input, 3/2
close.false %contradiction

%out = int.aff $expr, 7 {
    +3 @in.v[0].int
    -2 %integer_fact
    +1/2 %equality_fact
}
%out = round.hi $expr, %upper, %integer
%out = round.lo $expr, %lower, %integer
\end{lstlisting}

线性列表中的每一行恰好包含一个 multiplier 和一个 FactRef，中间至少一个空白，不加逗号。允许空列表；空 \code{lin.eq} 或 \code{lin.le} 只能证明零表达式。\code{int.aff} 的 offset 必须是精确整数。整数事实的 multiplier 必须为整数；EQ 项的 multiplier 可以是任意有理数，这是整数性通过等式传递的草案修订 P-04。

\code{drop} 的区间端点经符号解析后必须形成递增的连续 FactId 区间，其中每个 ID 均为当前 scope 已定义的 live fact。区间包含未来、dead 或其他 scope 的事实时不合法；不能因文本标签恰好具有连续数字而假定其机器 ID 连续。assembler 可以将 liveness 结果拆成多条合法 DROP 区间。

\subsection{Mapping 与 obligation 指令}

\begin{lstlisting}
map.copy @out.v[0], @in.v[1]
map.const @out.v[1], 0
map.aff @out.v[2], $expr

check.row.hi @out.r[0], %fact
check.row.lo @out.r[0], @in.r[1].lo
check.row.eq @out.r[0], %fact
check.var.hi @out.v[0], %fact
check.var.lo @out.v[0], @in.v[1].lo
check.var.int @out.v[0], %fact
check.obj.eq %fact
check.obj.le %fact
\end{lstlisting}

CHECK 的 proof operand 可以是任意合法 FactRef，包括静态模型事实和 branch assumption，无需先复制到 derived register。其类型和表达式必须与 kernel 生成的 obligation 精确一致。

\subsection{\code{exact} 简写}

\proposal{P-05}{为 CHECK 增加 \code{exact} 文本操作数；编译为 CHECK 的 EXACT flag，并由 kernel 直接检查下述有限规则。它不是 assembler 信任断言。}

\begin{lstlisting}
check.row.eq @out.r[0], exact
check.var.int @out.v[0], exact
check.obj.eq exact
\end{lstlisting}

exact 的判断对象始终是 kernel 自己由 output model 和当前 mapping 构造的表达式 $e$：

\begin{longtable}{p{0.18\textwidth}p{0.74\textwidth}}
\toprule
所需 FactKind & EXACT 可接受条件 \\
\midrule
\endhead
EQ & canonical $e$ 为恒等零。\\
LE & canonical $e$ 无变量项，且其常数不大于零。\\
INT & $e$ 的常数和所有系数均为整数；所有非零项涉及的 input variables 在 input model 中声明为整数或 binary。\\
\bottomrule
\end{longtable}

EXACT 不搜索 input rows，不自动匹配 input bounds，也不推导未知变量的整数性。相同 input row 直接用于 output row 时应显式写出对应静态 FactRef。一个必须依赖已验证 INT fact 的整数性检查，应使用该 fact 或 \code{int.aff}，不能要求 EXACT 发现依赖。

模型相关的 \code{exact} 预检查和 expression 归一化要求 assembler 获得两个 canonical models。即使 assembler 已预检查成功，verifier 仍必须独立执行同一判定；源文本不能成为可信输入。

\subsection{Structured cases}

\begin{lstlisting}
(%a, %b) = split.bin @in.v[0] {
case 0:
    yield %left_a, %left_b
case 1:
    yield %right_a, %right_b
}

(%a) = split.int.thr %integer, 13 {
case le:
    yield %left
case ge:
    yield %right
}

(%a) = split.int.range %integer, %lower, %upper, 3, 5 {
case 3:
    yield %value3
case 4:
    yield %value4
case 5:
    yield %value5
}

() = split.sign $expr {
case le:
    yield
case ge:
    yield
}
\end{lstlisting}

split 总是带输出 tuple；零个输出使用 \code{()}。\code{split.bin} 的 cases 必须按 \code{0}、\code{1} 顺序出现；\code{split.int.thr} 和 \code{split.sign} 的 cases 必须按 \code{le}、\code{ge} 顺序出现；\code{split.int.range} 的 cases 必须按闭区间中的整数递增排列。case labels 不是可跳转地址。

所有 sibling branches 必须从同一父状态开始执行。若一个 DEFER variable 在部分可达 branches 中已赋值、另一些可达 branches 中仍未赋值，merge 立即返回 \code{E_MAP}。只有所有可达 branches 均未赋值时才保留 UNASSIGNED；均已赋值时按 mapping merge 规则保留共同 expression 或封存为 SEALED\_PIECEWISE。

每个 case 必须以 \code{yield} 或 \code{close.false} 结束，之后不能有指令。\code{yield} 的 FactRef 数量必须等于 split 输出 tuple 长度；\code{yield} 不带操作数表示零个事实。所有可达 branches 在每个 yield 位置必须产生完全相同的 canonical proposition。已证明不可达的 branches 不需要 yield。split 可以嵌套，但不允许 JMP 或其他非结构化跳转。

\code{split.int.thr} 的 threshold 和 \code{split.int.range} 的两个端点必须为精确整数。对于 INT expression $g$，range 所用 lower fact 必须精确表示 $-g+L\le0$，upper fact 必须精确表示 $g-U\le0$。闭区间端点要求 $L\le U$。

对于 \code{split.sign}，两个 proof branches 的假设分别为 $e\le0$ 和 $-e\le0$，零点属于两支。mapping 的数值应用若遇到零，按草案 P-06 选择先出现的 \code{le} case；验证仍检查两个 proof branches。

\section{FactId、scope 与规范反汇编}

\proposal{P-03}{derived FactId 按语法顺序隐式分配：普通 producer 各分配一个 ID；split 先递归分配各 branch 的 producer IDs，再分配 merge outputs。ID 分配不依赖数学可达性。}

assembler 先解析全部符号和结构，依此分配 FactId；kernel 仍禁止未来事实引用。split 输出在源码头部命名，但只有 split 完成后可见，不能在其 branches 内引用。branch-local facts 不能通过符号泄漏到其他 sibling branch 或父 scope。源代码在文本上较早定义某个 fact，不表示该 fact 在另一个 scope 中可见。

规范反汇编将 expression 符号重命名为 \code{\$e0}、\code{\$e1} 等，将 derived facts 重命名为机器 FactId 的十进制形式。它输出固定 case 顺序、约分数值、展开后的 canonical expressions 和确定的空白规则。

\must{assemble--disassemble--assemble 必须保持 contract、数学 propositions、控制结构和 mapping 语义；基础 v0.1 不承诺逐字节还原原始文件。}

原因包括 scalar/expression pool 的合法重复条目、section 排列、debug 数据以及文本标签可能不同。若需要逐字节保存，工具必须提供独立的原始二进制导出模式，不能把这种需求隐藏在 canonical RPA 语义中。

\section{核心 EBNF}

以下 EBNF 使用 \code{[...]} 表示可选部分、\code{\{...\}} 表示重复、\code{|} 表示选择。\code{NL} 表示换行，\code{WS} 表示至少一个水平空白。注释在语法分析前处理；空白行允许出现在任意完整语句之间。表达式加减运算左结合；scalar 乘法优先于加减法，单目符号优先于加减法。

\begin{lstlisting}[caption={词法与顶层语法}]
digit       = "0" | "1" | ... | "9" ;
letter      = "A" | ... | "Z" | "a" | ... | "z" ;
hex-digit   = digit | "A" | ... | "F" | "a" | ... | "f" ;
WS          = (" " | TAB), { " " | TAB } ;
NL          = LF ;
uint        = digit, { digit } ;
sign        = "+" | "-" ;
integer     = [ sign ], uint ;
fraction    = [ sign ], uint, "/", uint ;
decimal     = uint, ".", [ uint ] | ".", uint ;
mantissa    = uint | decimal ;
scientific  = [ sign ], mantissa, ("e" | "E"), integer ;
number      = fraction | scientific | [ sign ], decimal | integer ;
identifier  = (letter | "_"), { letter | digit | "_" } ;
name        = identifier | uint ;
fact-name   = "%", name ;
expr-name   = "$", name ;
input-var   = "@in.v[", uint, "]" ;
output-var  = "@out.v[", uint, "]" ;
output-row  = "@out.r[", uint, "]" ;
fact-ref    = fact-name
            | "@in.r[", uint, "].", ("hi" | "lo" | "eq")
            | "@in.v[", uint, "].", ("hi" | "lo" | "int")
            | "@case", [ "[", uint, "]" ] ;
hash        = '"', 64 * hex-digit, '"' ;
objective   = "eq" | "le" | "none" ;
contract    = ".contract mip=", ("opt-equivalent"
              | "feasibility-equivalent"), NL,
              ".contract forward.obj=", objective, NL,
              ".contract reverse.obj=", objective, NL ;
file        = ".rpa 0.1", NL,
              ".source sha256", WS, hash, NL,
              ".target sha256", WS, hash, NL,
              contract, forward-domain, reverse-domain ;
forward-domain = ".forward", NL, domain-body ;
reverse-domain = ".reverse", NL, domain-body ;
domain-body = { expression-declaration },
              ".map", NL, { map-declaration }, ".endmap", NL,
              ".code", NL, code-body, ".endcode", [ NL ] ;
\end{lstlisting}

\begin{lstlisting}[caption={仿射表达式与静态 mapping 语法}]
expression-declaration = ".expr", WS, expr-name, "=", affine, NL ;
affine      = scaled, { ("+" | "-"), scaled } ;
scaled      = number, "*", signed-primary | signed-primary ;
signed-primary = [ sign ], primary ;
primary     = number | input-var | expr-name | "(", affine, ")" ;

map-declaration = "copy.range", WS, output-var, ",", input-var,
                    ",", uint, NL
                | "copy", WS, output-var, ",", input-var, NL
                | "const", WS, output-var, ",", number, NL
                | "aff", WS, output-var, ",", expr-name, NL
                | "defer", WS, output-var, NL
                | "defer.range", WS, output-var, ",", uint, NL ;
\end{lstlisting}

\begin{lstlisting}[caption={普通指令语法}]
linear-list = "{", NL, { number, WS, fact-ref, NL }, "}" ;
proof       = fact-ref | "exact" ;
instruction = "nop"
            | "drop", WS, fact-name, [ "..", fact-name ]
            | fact-name, "=", ("lin.le" | "lin.eq"), WS,
                expr-name, linear-list
            | fact-name, "=", "eq.pair", WS, expr-name, ",",
                fact-ref, ",", fact-ref
            | fact-name, "=", "box.le", WS, expr-name
            | fact-name, "=", "weaken.le", WS, expr-name, ",",
                fact-ref, ",", number
            | fact-name, "=", "int.aff", WS, expr-name, ",",
                number, linear-list
            | fact-name, "=", ("round.hi" | "round.lo"), WS,
                expr-name, ",", fact-ref, ",", fact-ref
            | "map.copy", WS, output-var, ",", input-var
            | "map.const", WS, output-var, ",", number
            | "map.aff", WS, output-var, ",", expr-name
            | ("check.row.hi" | "check.row.lo" | "check.row.eq"),
                WS, output-row, ",", proof
            | ("check.var.hi" | "check.var.lo" | "check.var.int"),
                WS, output-var, ",", proof
            | ("check.obj.eq" | "check.obj.le"), WS, proof
            | split ;
\end{lstlisting}

\begin{lstlisting}[caption={终止、split 与 case 语法}]
fact-list   = fact-ref, { ",", fact-ref } ;
output-list = fact-name, { ",", fact-name } ;
yield       = "yield", [ WS, fact-list ] ;
false       = "close.false", WS, fact-ref ;
branch-body = { instruction, NL }, (yield | false), NL ;
code-body   = { instruction, NL }, ("ret" | false), NL ;
split       = "(", [ output-list ], ")", "=", split-op ;
split-op    = "split.bin", WS, input-var, "{", NL,
                "case 0:", NL, branch-body,
                "case 1:", NL, branch-body, "}"
            | "split.int.thr", WS, fact-ref, ",", number, "{", NL,
                "case le:", NL, branch-body,
                "case ge:", NL, branch-body, "}"
            | "split.int.range", WS, fact-ref, ",", fact-ref,
                ",", fact-ref, ",", number, ",", number, "{", NL,
                range-cases, "}"
            | "split.sign", WS, expr-name, "{", NL,
                "case le:", NL, branch-body,
                "case ge:", NL, branch-body, "}" ;
range-cases = "case", WS, integer, ":", NL, branch-body,
              { "case", WS, integer, ":", NL, branch-body } ;
\end{lstlisting}

该 EBNF 规定结构；类型、范围、case 完整性、scope 可见性和终止可达性属于静态及执行语义检查。token 之间可插入水平空白，除显式 \code{WS} 外不强制数量。线性列表的右花括号后由外围语句换行结束。assembler 必须报告源码位置，不能通过自动补全缺失 case 或 obligation 使不完整程序成为有效证书。

\section{工具链接口草案}

\begin{longtable}{p{0.19\textwidth}p{0.73\textwidth}}
\toprule
工具 & 接口职责 \\
\midrule
\endhead
\code{rpa-as} & 解析文本，读取 canonical models 以检查索引和 exact 简写，展开 expressions，解析 FactId 和 scopes，输出 RPX。允许实现 liveness 优化，但不得改变命题或丢掉 obligation。\\
\code{rpa-dis} & 先验证 RPX 的结构，再输出 canonical RPA；不把成功反汇编解释为证明通过。\\
\code{rpa-objdump} & 显示 header、section directory、opcode counts、scope 深度及 fact liveness 等结构信息。结构统计与验证运行统计必须明确区分。\\
\code{rp-verify} & 读取两个 canonical models，绑定摘要和调用者要求，执行两个 proof tapes，成功后报告 VERIFIED 及 contract。\\
\code{mps-canon} & 将明确支持的 MPS 子集转换为 canonical model；不支持的输入特性必须显式拒绝。\\
\bottomrule
\end{longtable}

\begin{lstlisting}[caption={计划中的模型相关汇编与验证接口}]
mps-canon source.mps -o source.mcm
mps-canon target.mps -o target.mcm

rpa-as --source source.mcm --target target.mcm \
    proof.rpa -o proof.rpx

rp-verify --require opt-equivalent \
    source.mcm target.mcm proof.rpx

rpa-dis proof.rpx -o proof.canonical.rpa
rpa-objdump --sections --opcodes --liveness proof.rpx
\end{lstlisting}

以上命令展示接口契约而非可立即运行的实现。verifier 的成功输出至少应包含 VERIFIED、实际检查的 contract 和两个 model 摘要。instruction 数、derived fact 数、peak live facts 和耗时应说明其统计定义，特别区分 proof tape 解码、exact arithmetic 和模型加载成本。任何失败必须返回非零退出状态，并给出失败 domain、指令位置、错误类别和简短原因。

\section{扩展边界与保留语法}

基础 v0.1 不接受 \code{.batch} 或 \code{call.thm}。批量编码区 \code{0x80--0xBF} 和 \code{CALL\_THM} 的 \code{0xF0} 仅为未来扩展保留；基础 loader 遇到它们必须返回 \code{E_UNSUPPORTED}。

未来 batch 扩展必须定义完整的 core instruction 展开，包括 FactId 分配、scope、错误行为和 obligation 更新，不能通过批处理增加可证明命题。未来 theorem 扩展必须显式标识 trusted checker 和 TheoremId，不能用新名字绕过 trusted semantics 的版本审查。

\chapter{完整可复现实例}
\label{ch:examples}

三个 CLI 例子位于本手册的 \code{examples/} 子目录。每个目录包含完整 source.lp、target.lp
与 RPA 12 行式汇编 proof.rpa；汇编得到 RPX 11 字节，反汇编再汇编后与原文件逐字节相同。
\code{examples/python-api/} 另给出使用 Python 编译器 API 构造证明的脚本。

\section{等式消元}
源模型保留等式 $x_0+x_1=1$，目标只使用 $y=x_1$；反向恢复 $x_0=1-y$。
\rpxlisting[caption={完整源模型}]{examples/elimination/source.lp}
\rpxlisting[caption={完整目标模型}]{examples/elimination/target.lp}
\rpxlisting[basicstyle=\footnotesize\ttfamily,caption={两个方向及全部位置义务}]{examples/elimination/proof.rpa}

\section{冗余行与分段恢复}
\code{examples/redundant/} 通过已有界证明冗余行，同时两个方向保留完整目标。
\code{examples/piecewise/} 使用穷尽分支实现分段恢复；每支仍须独立完成全部输出义务。
完整源文件是本手册组成部分，可直接阅读、汇编和验证。

\section{Python 编译器 API}
\code{rpx} 的 Rust Core 持有模型空间、作用域、事实、映射和完成状态。
作者选择推导；编译器不调用 solver，不搜索证明。ModelBuilder 的 build 封闭建模阶段，
finish 封闭证明叶子。不同模型的表达式不能按同名变量自动混用。

\begin{lstlisting}
from rpx import Model, Proof
source = Model.read("source.lp")
target = Model.read("target.lp")
proof = Proof(source, target)
f = proof.forward.map({"y": source.vars["x1"]}, default="same_name")
f.prove("objective.non_increasing", by=-f.row("balance"))
f.finish()
proof.reverse.map({"x0": 1-target.vars["y"],
                   "x1": target.vars["y"]}).finish()
bundle = proof.compile("build/elimination", target_artifact="target.lp")
print(bundle.verify().raise_for_status())
\end{lstlisting}

\code{map} 的默认名称匹配在 Rust 中展开。变量、行、目标、trace 和 certificate 是
原生视图，整体物化必须显式请求。Fact 算术产生线性组合；它不与普通表达式混同。
连续 split 返回两个 child scope，分别带有 $e\le a$ 和 $e\ge a$ 的假设，边界属于两支。
整数 split 要求整性证明，其右支是 $e\ge a+1$。Cases 的左右推导 scope 分别提交共同
结论或矛盾，然后向父级返回一个事实。Interval、Menu、Signed、Flow 和 Subflow 使用
类型化输入与生成句柄。父空间可嵌入扩展，生成变量不能逃到父级或兄弟扩展。

封闭证明只 lowering 一次。RPA disassembler 与 RPX encoder 直接消费同一份原生
Certificate。Bundle 保存 source、target、RPA、RPX 和目标原文件；manifest 使用包内
相对路径与内容摘要，可搬移。Chain 复制阶段包并校验连接。每次验证使用本次预算，
历史 receipt 不替代重新验证。Kernel 结论和 artifact audit 分别显示。

五个可执行脚本与 \code{tutorial.ipynb} 位于 \code{examples/python-api/}。
\code{elimination.py} 展示建模和消元，\code{redundant.py} 展示冗余行，
\code{piecewise.py} 展示连续 split，\code{interval\_lift.py} 展示生成变量，
\code{load\_models.py} 展示文件读取和稀疏映射。脚本默认使用临时交付目录。
从仓库根执行：
\begin{lstlisting}
make -C docs/rpx check-compiler-api
\end{lstlisting}

\section{统一检查命令}
从仓库根执行：
\begin{lstlisting}
python3 tools/check_rpx_manual.py --kernel target/release/rpx
\end{lstlisting}
当前检查执行 39 个 CLI 命令：验证三个完整例子、汇编以及反汇编/再汇编字节往返，
并运行 12 个负例检查，覆盖缺失义务、旧版本头、错误反向映射和资源耗尽。
三个例子均各有 9 项总 Return 义务。记录保存在 \code{examples/verification.json}。
汇编阶段拒绝旧版本头/语法错误与执行阶段拒绝数学错误属于不同检查边界。

\chapter{执行成本、资源预算与测量}
\label{ch:performance}

本章描述当前 \code{crates/core/src/verify.rs}、\code{api.rs} 的计数与实现，以及仓库已保存的受控实验。它不把 D0 的 streaming tape、8-byte LiveFact、mmap 或 batch 设想当作现有性能。当前 CLI 读取 owned bytes，decoder 构建完整 owned model/certificate，再执行 verifier；只读输入、已解码对象、精确算术临时量和 proof scopes 均会占内存。

\section{公开计数的确切含义}

\code{VERIFIED} report 的 \code{statistics} 由 kernel 生成，不是作者提交的证书字段。下面列出主要计数；全部字段定义在 \code{Statistics}。

\begin{longtable}{p{0.31\textwidth}p{0.61\textwidth}}
\toprule
字段 & 当前计数边界 \\
\midrule
\endhead
\code{obligations} & 实际执行的 terminal/output checks，包括方向 Return、严格改进和 injectivity 等通过同一 obligation 入口的检查。Split 的两支均计数；不含 lift premises、local Cases Witness、point checks 或 fiber 内部检查。\\
\code{linear_terms} & Lin 展开所访问 premise expressions 的非零 coefficient 项数，不是 Lin 数或 scalar arithmetic 次数。\\
\code{substitution_terms} & substitute 展开每个被引用 map expression 时访问的非零项总数；不同 rows 的重复访问重复计数。\\
\code{branches} & ordinary Split 与 local Cases 的 branch points；一个二支 split 计一个 point，而非两个 leaf。\\
\code{peak_live_facts} & 同一时刻当前 scope 可见的 logical local-register bindings 的峰值；不是累计定义数，也不是物理存储或 RSS。\\
\code{work} & 被明确 charge 的操作总量，包括 step、expression/fact/map 项及常数、ancestor lookup probes、构造器和相关验证工作。是资源政策单位，不是 CPU instructions、运行时间或严格 bit-complexity 上界。\\
\code{points}, \code{point_coefficient_terms} & 完整 point checks 及 point/observation 精确 evaluation 项数；不含 fiber 反证中的 Lin/代换项。Point 的 bounds/rows 检查不进入方向 Return 的 obligations。\\
\code{plotkin_pairs} & Plotkin witness 中已检查的 coordinate pairs；还存在 labels、membership 与 premise 扫描成本。\\
\code{flow_*}, \code{interval_*}, \code{menu_*}, \code{signed_*}, \code{subflow_*} & 各类 lift 次数及其 generated variables、rows、terms。记录完整 generated extension，不代表存在性引理已被另一个 universal checker 证明。\\
\code{bijections}, \code{injectivity_*} & V7 paired-domain bijection 与其完整 injectivity domain/checks。Legacy \code{inverse} 的成功不自动增加此 bijections 计数。\\
\code{objective_level_*}, \code{halfspace_*}, \code{improvements} & 条件域构造次数、生成行和项，以及严格改进 witnesses；V10 的新增 zero counters 会从 JSON 中省略。\\
\bottomrule
\end{longtable}

Context 的 \code{charge} 做 checked addition，超过有效 work ceiling 就返回 \code{RESOURCE_LIMIT}。例如一次 static/local fact 读取还计入 expression 的 terms+1；local 引用另外计 ancestor 查找 probes。一次 map substitution 按引用表达式的 terms+1 charge。换一种内部计数政策可能改变 work，必须在实验中说明；同一 work 不能保证同一机器上耗时相同。

第\ref{ch:examples}章的分段证书只有一组三项 Forward 数学要求，但 verifier 在两个 Return 各检查一次，所以 Forward 的 obligations 是六。完整 fiber refutation 使用同一 verifier 的 fresh-scope 调用，但 \code{point.rs} 只把其 work 加回共享预算，不把内部 branches、linear\_terms、substitution\_terms、peak\_live\_facts、plotkin\_pairs 或 obligations 汇入顶层 Statistics；分析这类 proof 时不能从顶层零计数断言内部没有该工作。无限精度 rationals 的位长不同，也会使同样项数的 Lin 成本很不同。不能把任何一个单项计数当作完整成本。

\section{共享预算与可复核的资源结果}

当前默认 \code{Limits} 为 work $10^9$、live facts $10^6$、depth 64，root depth 为 1。一个 request 的全部方向、条件域、lift、injectivity、points/fiber refutations 和 improvements 共用 work；嵌套 scope 或新方向不会重置。CLI 可配置 work 与 live facts，depth 不提供 CLI override。

\begin{lstlisting}[caption={预算是实际 verifier 参数}]
target/release/rpx verify --work 1000000000 --live-facts 1000000 cases/test58/source.mcm cases/test58/target.mcm cases/test58/proof.rpx cases/test58/request.json
\end{lstlisting}

有限预算是部署政策，不是模型数学定义。\code{INVALID}、\code{RESOURCE_LIMIT} 和 \code{IO_ERROR} 都不建立 reformulation 不成立。第\ref{ch:examples}章已在同一正例上实际观察 \code{--work 0} 的 RESOURCE\_LIMIT；test36 的同一完整产物在 $10^9$ work 政策下成功，旧 $10^8$ ceiling 下耗尽预算。这不是数学关系改变，也不能用失败运行时间作 speedup 的分母。

还必须分别记录其他层次的边界：CLI 每文件 owned read 为 512 MiB，RPA/request 文本为 256 MiB，serde JSON nesting capacity 为 128，binary nesting 为 64；模型 dimensions、counted arrays 和 rational magnitude 有各自上限。proof depth 与 JSON wrappers 的嵌套数不同。当前精确算术中间 numerator/denominator 的 bit budget 为 65,536。提高 work 不会绕过这些独立边界。

\section{代换、精确累加与表示成本}

对 output row $r$ 和完整 map $y_j=E_j(x)$，直接代换展开
\[
A_rE(x)-u_r=c_r+\sum_{j\in\operatorname{supp}(A_r)}a_{rj}E_j(x),
\qquad
W_r=\sum_{j\in\operatorname{supp}(A_r)}\operatorname{nnz}(E_j).
\]
输入项可能重合、相消，输出可能变稠密；$W_r$ 不只由 output row 的 nnz 决定。跨 row 重复的代换仍付出访问成本。当前 \code{Accumulator} 用 hash map 暂存精确 coefficient，再排序生成 canonical sparse expression；hash 只选择容器 bucket，精确 key/rational equality 决定接受。没有通用 affine-expression DAG、共享 mapping DAG 或证书中间结果缓存的格式承诺。

\code{BigRational} 保持精确数学语义。现有 denominator-one、乘 $\pm1$ 和零加值快路径避免部分通用算术工作，同时保留所需 intermediate/result checks；非 unit rational 仍使用 arbitrary-precision 运算。它不等于 D0 建议的机器小整数与大整数自动提升 ABI，也不把 coefficient 位长或 gcd 成本计成 CPU 指令。

展开长 prefix maps 可有二次物理大小；scalar injectivity proofs 也可能很大。二进制减少 JSON 文本开销，并未保证完整 proof resident memory 与 peak live facts 成正比。一个解释性的成本分解是
\[
T=T_{\rm read}+T_{\rm decode}+T_{\rm execute}+T_{\rm process},
\qquad
T_{\rm execute}=T_{\rm algebra}+T_{\rm scopes}+T_{\rm builders}+T_{\rm checks}.
\]
它不是已证明的统一渐近上界；精确算术位长与保存的 physical representation 需要额外描述。

\section{Shared scopes 的实现与内存边界}

当前 \code{Facts} 借用不可变 parent，只在 child 保存 local definitions 与 deleted-ID masks。其可见集合为
\[
L\cup(P\setminus D),
\]
其中 local bindings 优先；drop/reuse/drop 不会复活祖先旧值。一个 child 的 DROP 不改变 sibling 或 parent。Cases 的两支独立检查，仅把共同结论插入 parent；ordinary Split 没有后续 continuation，left 检查后 right 可接管未改变的 parent 状态。这消除了每个 branch 复制全部 inherited facts 的成本。

祖先 lookup 的 probes 被计费，深度仍有界；local insertion 和 deletion-mask growth 使用 fallible reservation。\code{peak_live_facts} 只计当前可见 bindings。被 child mask 隐藏的 immutable ancestor payload 仍在内存中，masks 也占空间，所以 logical live limit 不能直接转换成固定 bytes/live-fact 或 peak RSS 上限。decoder 的 owned objects、arithmetic allocations 及部分依赖库 allocation failure 也是当前内存保证的边界。

独立 copied-BTreeMap state oracle 的 127,072 次操作、depth-64 checks 与 negative scope tests 见 \code{audit/SHARED_SCOPE_REVIEW.md}；生产集成另见 \code{audit/scope_v6_merge.json}。状态正确性审阅和 latency 测量分别保存，不能互相代替。

\section{时钟字段与端到端成本}

当前 Rust façade 和 CLI 使用同一报告函数，公开 microseconds 字段：
\begin{longtable}{p{0.28\textwidth}p{0.64\textwidth}}
\toprule
字段 & 含义 \\
\midrule
\endhead
\code{read_us} & 读取 source、target、RPX 和 request；filesystem-free \code{verify_bytes} 不返回此字段。\\
\code{model_decode_us} & 两份完整 MCM 的 decoding/validation 阶段。\\
\code{proof_decode_us} & RPX decoding 加 caller request text parsing。\\
\code{execute_us} & 完整 \code{verify_request_with_diagnostics} 调用，包含全部请求验证。\\
\code{total_us} & kernel API 本次调用的内部总时间，包括报告前的操作；不含 process startup/shutdown。\\
process wall time & 外部计时，另含进程启动和退出。不是另一个 kernel 公开计数字段。\\
\bottomrule
\end{longtable}

MPS/LP import、作者生成 RPA、assembler 和任何外部 oracle 必须单列。\code{CANONICALIZED}/\code{ASSEMBLED} 的 \code{total_us} 不能当作 verification time。未重测的历史 preparation 应标为 archived/unmeasured；raw 数据中的空 preparation array 或 sentinel 0 不能解释为零成本。

可靠的对比先固定完整数学 request、模型/proof bytes、source/binary/lockfile identities、硬件、编译选项和 budget，再记录 cold/warm 文件状态、重复次数、顺序与 isolation。每个样本都核对退出码和 VERIFIED，保存失败和 outliers；timing 与 RSS 另起 processes。至少报告 min/median/max 或完整分布，并同时保存计数、模型 dimensions、artifact size 与阶段时间。

\section{已有实验：明确保留适用范围}

\subsection{共享状态与精确算术}

\code{bench/SHARED_SCOPE_REPORT.md} 固定了 frozen V5 的 before/after sources，数学规则和 certificate bytes 相同，仅改变 scope 表示及相应 work accounting。三个 copying-heavy workloads 的 execution medians 如下，单位 ms：
\begin{center}
\begin{tabular}{lrrrr}
\toprule
active facts / Cases & before & after & work before & work after \\
\midrule
256 / 256 & 3.439 & 0.425 & 267,535 & 8,720 \\
1,024 / 1,024 & 51.413 & 1.605 & 4,215,823 & 34,832 \\
4,096 / 4,096 & 827.142 & 6.517 & 67,194,895 & 139,280 \\
\bottomrule
\end{tabular}
\end{center}

这支持消除所选二次复制成本，不支持普遍 127 倍加速。完整 broader matrix 保存了 nonbranching regressions：例如 real65 execution median 从 3690.334 增到 3793.140 ms；logical counters 中除了 work 以外的数学项、全部 physical artifacts 一致。test18 后来的成功是原字节的预算内 replay，其原失败没有成功性能分母。

\code{bench/RATIONAL_SHORTCUTS_REPORT.md} 对 49 个 paired workloads、1,176 个成功 verifier processes 观察同 bytes、同完整计数下的算术快路径。所有 execution medians 降低，但 startup 主导的 process latency 不全部降低；例如 injectivity-64 从 2.186 增到 2.391 ms。RSS median 变化在 3\% 内。报告描述其固定 prototype binary，不能把后来 resource/API 修订后的可执行文件自动等同于被测候选。生产合并边界与额外检查记录在 \code{audit/V7_ARITHMETIC_MERGE.md}。

\subsection{V10 旧任务 regression 与顺序敏感性}

\code{bench/V10_REGRESSION_REPORT.md} 保留 31 个 paired 旧任务、712 个成功 verification/RSS processes，相同 artifacts、request 和完整 non-timing results。candidate 在 19 个任务的 execution median、23 个任务的 full-process median 上增加。test65 execution 为 1792.479 到 1864.392 ms，约 $+4.0\%$；rational-1000 为 0.792 到 0.972 ms，增加 0.180 ms。没有通用加速或无 regression 声明。

预先声明的一次 alternating-order follow-up 保存另外 72 个 timed calls 和八个 excluded warmups，四个任务的 paired median difference 分别为 rational-1000 $+0.003$ ms、test65 $+82.552$ ms、test71 $-0.038$ ms、test79 $-0.005$ ms。test65 的正成本保留为交付限制，其他三项原有较大差异未在本轮持续。分布和 within-pair order effects 不支持把一个精确百分比归因于代码。原矩阵未被替换，也没有新的 RSS 声明；原始报告见 \code{bench/results/v10-order-sensitivity/ORDER_SENSITIVITY_REPORT.md}。

\subsection{新能力绝对 baseline}

\code{bench/NEXT_CAPABILITY_DRAFT_BASELINE.md} 固定 22 个任务，264 个成功 verification/RSS processes，并单列 63 次 synthetic author/assembly preparations。使用最终 cold checkpoint 的同一 binary，task-owned build/preparation/acceptance 工作在采样期间 idle；warm files、fresh process，OS activity 未控制。新任务不与无法证明相同命题的旧 binary 计算 speed ratio。

两个已有报告的 median 样本如下；完整 min/median/max、decoder 时间和身份都在原报告：
\begin{center}
\begin{tabular}{lrrr}
\toprule
任务 & execute ms & process ms & peak RSS MiB \\
\midrule
完整 actual test58 & 18.369 & 28.749 & 31.891 \\
subflow 1,024 nodes / 4 groups & 56.732 & 87.707 & 69.141 \\
\bottomrule
\end{tabular}
\end{center}

第二项的完整 Python preparation median 为 9.154 s，包含简单 $O(nm)$ author construction；不能把这个成本藏在 constant-time theorem 的叙述中。actual test58 的历史 author preparation 未重测，表中没有零成本含义。Corpus 的总 functional replay 耗时和本手册单次示例 replay 都不是该受控矩阵的替代。

\section{复现与下一次优化的登记}

仓库已有 fixed matrix drivers；新的运行须用未使用的 label 保存 raw evidence。下面的命令可建立当前 binary 的新绝对测量，不会复写历史样本：
\begin{lstlisting}[caption={当前完整初始与 extended matrices}]
python3 tools/benchmark.py --kernel target/release/rpx --source-root . --label manual-current-initial
python3 tools/benchmark_extended.py --kernel target/release/rpx --source-root . --label manual-current-extended
\end{lstlisting}

正式执行前应安排独立采样窗口并登记机器与其它任务状态；这些命令不自动保证 isolation。固定 V10 final-driver 与能力 config 还验证专门 binary/manifest hashes，因此重建当前 binary 不会被默认为同一历史样本。已有 driver 拒绝覆盖 raw evidence，应换新 label 或单独建立新固定计划。

共享 DAG、dedicated batch、streaming decoding、模型缓存和独立 solution reconstruction 是可能的后续工作，均非当前报告中的既有收益。新增 batch 必须另定义可观察语义、operand 检查、scope/obligation 与失败条件，不能用名称加入隐藏 theorem。优化接受路径继续做精确验证；浮点 residual、probabilistic hash 和外部 solver success 不构成接受条件。

任何对其它 proof format 的速度比较，都须提供相同数学请求和推导、相同线程政策及 preparation/verification 边界。当前报告没有这种受控 cross-format 实验，因而不宣称相对 VIPR 或求解器的性能优势。

历史 source snapshots 与部分 prototype 完整副本已在交付后的清理中删除，
原 manifests、receipts、logs 与 diffs 保留；具体清理边界见
\code{audit/POST_DELIVERY_CHANGES.md}。依赖这些副本的历史重建必须先恢复原版本数据。
本章的新测量入口使用当前源码，不能把重新编译的当前 binary 当作原实验 binary。

\appendix
\chapter{设计沿革、冻结范围与后续登记}
\label{app:revisions}

本手册以当前已实现 kernel 为对象。原 D0 是 2026 年 10 月 9 日整理的设计输入，其 \code{.rpa 0.1}、28-opcode tape、128-byte 候选 header、digest 绑定、FactId SSA 和分支汇合不是已实现格式。MCM 1、RPA/RPX 1--10 均已冻结；V9/V10 的明确冻结决定于 2026-10-14 登记，保持现有字节与数学条件。D0 候选 ABI 仍仅为历史设计输入。

\section{D0 提案 P-01 至 P-07 的实际处置}

保留原提案编号以便追溯。以下为当前实现与原建议的对应关系，既非待表决草案，也不表示原字段编码已经发布。

\begin{longtable}{p{0.09\textwidth}p{0.18\textwidth}p{0.65\textwidth}}
\toprule
编号 & 处置 & 已实现行为与差异 \\
\midrule
\endhead
P-01 & 数学原则落地，接口重写 & 目标统一为完整 affine minimization，包括常数。caller request 独立且必须与证书完全相等；MI/LP 用独立 input/output integer booleans 表达，不是 D0 两个 profile。双向 Equal 保留 attainable objective values 与 extended infimum；双向 NonIncreasing 保持 extended infimum 和非空性。映射无需互逆；bijection/projection 是另行绑定并核验的要求。\\
P-02 & 用完整 terminal maps 替代 & 每支从同一 parent 的 isolated fact view 开始。每个 Return 一次给出全部 output coordinates，核验全部 obligations；Split 是 terminal，没有 post-merge continuation。无 DONE 交集、DEFER、三态 map 或 SEALED\_PIECEWISE。V2 Cases 仅导出共同的已核验 proposition，不能导出独立 existential output fragments。\\
P-03 & 不采用 D0 SSA 编号 & Local 是 mutable u32 register。定义不能覆盖 live binding；DROP 后同一 ID 可重新定义，引用立即解析。没有按语法深度优先预分配的历史 FactId slots。child 的 local writes/deletions 隔离，borrowed ancestors 与 deletion masks 保留 state invariant。\\
P-04 & 数学闭包落地于 Lin & \code{Lin} 的 Int kind 允许有理数倍 EQ premises；constant 和 INT multipliers 必须整数。EQ=0 的项不改变整数性。启用 input integrality 才能读取 declared integer predicates。由模型 facts 的推导与 Exact 的结构整数性判定不同。\\
P-05 & 精确简写落地，编码不同 & \code{Exact} 是 terminal/derivation witness Ref，kernel 自行检查 constant-zero EQ、constant-nonpositive LE 或有整数 domains 支持的 integral affine expression。不是 CHECK flag，不产生 local fact，也不能作为 Lin premise。\\
P-06 & 覆盖原则落地，重建接口未发布 & Real Split/Cases 的 closed guards 为 $e\le0$ 与 $-e\le0$；integer guards 为 $e\le0$ 与 $1-e\le0$，须有匹配 INT。两个完整子树都被 decoder/verifier 检查。D0 的独立数值解应用与零点 le 优先 CLI 未实现；其建议没有成为接受规则。\\
P-07 & 候选 wire layout 被替代 & 当前采用 16-byte common header、little-endian length/count 与嵌套 records，canonical rationals 和 complete model/claim objects。无 ScalarRef immediate、pool/directory tape 或可选 metadata section。实际 tags、version gates、reserved bytes、preflights 见第\ref{ch:format}章；原候选字节向量不能作当前 golden vector。\\
\bottomrule
\end{longtable}

这些处置保留目标：独立模型绑定、精确有限证书、完整双向语义、可审阅失败与资源边界。实际模型绑定依赖每个 input premise 与完整 output obligations 的语义核验，RPX 不含允许跳过模型检查的摘要。准备脚本和 manifest hash 是来源证明，不是数学公理。

\section{实现版本 V1 至 V10 的能力历史}

各版本的数学与 byte meanings 保持；十份 FORMAT increment 已去重合为 \code{docs/FORMAT.md}。表中的版本是使用相应 feature 的最低 RPX 版本；assembler 通常输出实际必要的最低版本，RPA N 标头仅允许 N 范围的 features。不是每个版本单独运行一个 verifier。

\begin{longtable}{p{0.09\textwidth}p{0.14\textwidth}p{0.69\textwidth}}
\toprule
版本 & 协议状态 & 新增能力与边界 \\
\midrule
\endhead
V1 & 已冻结 & 完整模型 affine directions；Lin/Round/Drop，完整 Return、exhaustive Split、Absurd；独立 caller claims，objective/projection/legacy affine inverse。MCM 1 精确模型与严格 JSON RPA 1。来源：\code{SEMANTICS.md}、\code{FORMAT.md}、\code{RPA.md}。\\
V2 & 已冻结 & EqPair 的精确 antisymmetry，local Cases 的共同 proposition derivations；保留全部 simultaneous output obligations，避免重复无关整模型分支。实际 test88 integer alias 和 test84 parity 见 \code{FORMAT.md}。\\
V3 & 已冻结 & 静态 Plotkin theorem、caller objective lower bound、完整 Feasible 和 full-permutation Strictness points。test71 核验 integer 等价、严格 LP strengthening 和精确 source-integer/target-LP 最优值。静态 theorem 无 host callback。来源：\code{FORMAT.md}、\code{POINT_V3.md}。\\
V4 & 已冻结 & Common-source joint FlowLift：所有 commodities 一起满足每个节点 balance，并保留每条弧总流及 circulation。test28 不把单独局部存在性当作 joint tuple。来源：\code{FORMAT.md}、\code{FLOW_V4.md}。\\
V5 & 已冻结 & Independent IntervalLift 与 caller affine observations；完整 old-scope premise checks 后联合生成所需连续坐标。test23 使用 signed affine observations。来源：\code{FORMAT.md}、\code{INTERVAL_V5.md}。\\
V6 & 已冻结 & FiberStrictness 和完整 point refutation payload；固定全部 caller coordinate/affine observations，在实际完整 input LP fresh scope 中反证整个 fiber。Shared-scope production integration 是表示/资源优化，非新数学 theorem。test18/test19 见 \code{FORMAT.md}、\code{FIBER_POINT_V6.md}。\\
V7 & 已冻结 & Optional MenuLift 和 explicit complete-domain bijection pairs；forward-after-reverse identities 与 paired-domain injectivity 覆盖全部坐标。test59/test60 核验整数 menu 双射与 LP observation images。来源：\code{FORMAT.md}、\code{MENU_V7.md}。\\
V8 & 已冻结 & Activated signed real decomposition，同时生成 $p,n,u,v$，证明 $p-n=F$、$u+v=Z$ 和容量 bounds。必须 integer activation，而 real flow 可以连续。来源：\code{FORMAT.md}。\\
V9 & 已冻结 & Caller-bound \code{input_objective_level} 把完整 input domain 与完整 objective equality 相交；legacy opposing inverse 要额外 image-level membership。条件 bound 不能充当 global bound；受限 claim 不偷换旧 whole-domain point/bijection 的含义。test80 与 \code{FORMAT.md}。\\
V10 & 已冻结 & 严格 positive-domain improvement、完整 affine ImageExclusion、joint demand-bounded SubflowLift，以及 caller-bound original-coordinate \code{input_halfspace}。追加 records 与 feature flags 独立 gating；兼容已有有效字节。实际 test36/test58/test62/test83，来源：\code{FORMAT.md}、\code{SUBFLOW_V10.md}。\\
\bottomrule
\end{longtable}

V1--V8 的 checkpoint test/case counts 是历史时点的证据，不应改写成今天的总量。例：frozen V8 rebuild 是 198 tests、当时列举的 34 个实际 certificate；后来的交付 cold build 是 262 tests。两个记录各有 source/binary/lock manifest，不能用后者倒填前者。统一的 \code{docs/FORMAT.md} 覆盖 V1--V10；此次去重不创造新版本，V9/V10 另由明确项目决定冻结。被旧 manifest 固定的原规范字节移入 \code{audit/specification-history/}，历史路径和哈希含义保留。

当前 façade 的 \code{api.rs} 与 CLI 共用 report/status/exit handling，提供文件及 filesystem-free verification、canonicalization、assembly、inspection 和 model export。这是交付 snapshot 后的 interface 整理，未改 wire formats 或数学规则。精确差异在 \code{audit/POST_DELIVERY_CHANGES.md}；当前机器可读 \code{code}/\code{location} 为诊断，不给 proof 更多接受权。

\section{D0 开放项 O-01 至 O-05 的当前状态}

\begin{longtable}{p{0.09\textwidth}p{0.18\textwidth}p{0.65\textwidth}}
\toprule
编号 & 当前状态 & 已解决部分与保留边界 \\
\midrule
\endhead
O-01 & MCM 1 已实现并冻结 & 独立 canonical model 文件、rational/name/dimension validation 与 encoder/decoder 已公开。原 SHA 覆盖范围问题不再是接受前提，因为 proof 采用全模型语义绑定。\code{export-model} 输出加载后的全部 canonical JSON，不是 alternate acceptance representation。\\
O-02 & 已定义可信输入子集 & \code{MPS_V1.md} 与 \code{LP_V1.md} 明确 defaults、bounds/ranges、names、exact decimal、objective sense 和 unsupported constructs。全部 159 个 corpus MPS 与三个 LP 已导入；LP 与独立 Fraction parser 比对全部字段。支持这个子集不承诺所有厂商 dialect、SOS、indicator、quadratic 或 semicontinuous semantics。\\
O-03 & 独立数值解工具未发布 & Return 与 piecewise proof 的数学 total-map 语义已核验，完整 inverse/bijection 的要求已实现。当前 CLI 没有统一 numeric solution application/export 命令；不能把 verify、disassemble 或 model export 当作该工具。新的应用接口需另定义 guards、扩展 witness 实例化与 exact output。\\
O-04 & 有限静态扩展已实现，通用 batch 未实现 & V3--V10 通过 reviewed static constructors/checkers 扩展能力，未知 tags 与不支持的版本拒绝。无 theorem registry、动态插件、case whitelist 或 host callback。通用 expression/map DAG、dedicated batches、capability negotiation 仍待独立设计和评审。\\
O-05 & 有明确当前政策与真实 benchmarks & 当前预算、CLI overrides、四类 verifier outcomes、固定矩阵和完整 raw evidence 已落地。默认 work 从旧 $10^8$ 提到 $10^9$ 是有记录的 resource-policy 变更，不是速度优化。严格 bit-complexity/peak-RSS 上界、全部 allocation failure 可恢复性和所有未来任务均能预算内验证没有被证明。\\
\bottomrule
\end{longtable}

O-01/O-02 已不阻止当前工具链互通；O-03/O-04 的未来 representation/application 工具也不是当前 81 个实际任务完整验收的前提。资源政策与数学 semantics 仍需区分：预算耗尽不使完整证书成为逻辑反证。

\section{V9/V10 冻结决定与兼容边界}

项目于 2026-10-14 明确冻结 V9/V10，基线与 fresh locked tests、完整实际语料回放见 \code{audit/PROTOCOL_V9_V10_FREEZE.md}。冻结保留以下语义和兼容边界：
\begin{enumerate}
\item V9 objective-level 与 V10 halfspace 原始坐标域、相交次序、生成行顺序及 opposing legacy inverse membership；受限 claim 与旧 point/bijection 的禁止组合。
\item PositiveRefutation 的非严格 $P\le0$ 结论与旧 scope $P\ge0$ assumption、ImproveReturn 的正 margin，以及所有 positive-domain points 的完整 simultaneous self-model map。
\item ImageExclusion 的完整 caller-bound map、完整实际 LP 点、actual bound/row violation；不得宣传为未检查的 global containment。
\item Subflow 的 demand-bounded joint-existence 引理、每一组所有 old-scope premises、完整 generated schema 和 output tuple；builder 仍是可信 inference code。
\item 每个新增 record 的版本 gating、flags/reserved bytes、count/remaining-payload preflights、depth/work/live limits；所有 syntactic branches 均受检查，包括可能矛盾的 branches。
\item 既有 V1--V8 和早期 V9/V10 有效字节含义与 absent-field bytes 保持；历史 78 requests 的兼容 receipts 保留，不能以重建文本替代原字节对照。
\item 完整 actual-case caller scope：成功提交的弱 request 与 original-report 完整要求必须分开复核，不能以模型/proof hashes 或 author scope hints 掩盖遗漏的 property。
\item Fresh locked source build、完整 tests、cold actual-model replay 与既有独立 scope review、受控成本矩阵分别有身份记录；冻结决定不扩展测试证据为普遍性能承诺。
\end{enumerate}

最终交付已经完成 81/81 完整实际 cases、11 个 families、83 次成功 kernel calls；test46 三个独立 direction files 的 conjunction 在 delivery metadata 中明确登记，没有 cross-file facts/lemmas。原 78 个 submitted requests 也维持 byte/non-timing compatibility。这些已完成事项见 \code{audit/FINAL_DELIVERY_REPORT.md}、\code{audit/acceptance.json}、\code{audit/contract_scope_gates.json}。旧 \code{DELIVERY_ACCEPTANCE.md} 的 76/81 与个别历史 README 的 pending 文句记录此前 scope 审查，不能取代最新 final ledger。

可信边界包含输入与证书解析、精确算术库、表达式与证明检查器、静态扩展 builder，
以及编译器、运行时和操作系统。独立 oracle、数学论证、攻击测试与完整语料互为补充，
不构成 Rust 实现的形式化证明。

\section{手册修订记录与维护准则}

\begin{longtable}{p{0.16\textwidth}p{0.17\textwidth}p{0.59\textwidth}}
\toprule
日期 & 修订 & 内容 \\
\midrule
\endhead
2026-10-09 & D0 provenance & 整理原 28-opcode v0.1 候选设计、P-01--P-07、候选 wire layout 与 placeholder-hash 示例。它是设计来源，没有成为已实现 ABI。\\
2026-10-13 & 当前实现手册 & 按 safe Rust kernel 重写模型/请求、scalar/local/theorem rules、joint maps、MCM/RPX records、严格 JSON RPA、CLI/API、预算与诊断；完整区分 frozen V1--V8 和 working V9/V10。\\
2026-10-13 & 可执行证据 & 新增三份完整双向 LP/RPA/request，实跑三正例及回读、九次 INVALID 逻辑变异、三次 RESOURCE\_LIMIT；另复核完整 test36/test58/test62/test71/test80，保存实际产物与报告身份。\\
2026-10-14 & 协议冻结与规范合并 & 明确冻结 V9/V10；十份 FORMAT 与重复的 CURRENT reference 统一为 FORMAT.md，保存原始 hash-pinned 规范；同步中文手册、引用及完整验证基线。\\
2026-10-14 & 当前文档清理 & 移除已实现能力的 Flow、Interval、Menu 与 objective-level 旧提案；语义与 RPA 教程覆盖 V1--V10，修正公开 Rust 接口说明并同步引用。\\
\bottomrule
\end{longtable}

维护本手册时，先检查当前 code 的 public types、decoder/version gates 和实际 checker，再同步 request grammar、全部义务、例子与失败分类。性能优化要保存固定任务和原始数据，数学扩展要给明确规则与可信边界，冻结状态要有独立决定。更新中文手册不回写历史 checkpoint 的证据，也不从完整 corpus 的成功推导所有可能 MILP 改写都有紧凑可接受证书。

\end{document}
